% !TeX program = xelatex
% 数学 LaTeX 排版 Skill v4.0.0；版式保持 v3 金标准. XeLaTeX 编译至少两次. 不分发字体.
% WZ-LOCKED-CLASS-BEGIN
\begin{filecontents*}[overwrite]{elegantbook-original-adapter.cls}
\NeedsTeXFormat{LaTeX2e}
\ProvidesClass{elegantbook-original-adapter}[2026/09/23 v3.0 fixed mathematical book environments]
\LoadClass[10pt,letterpaper,oneside]{book}
\RequirePackage[UTF8,scheme=plain,fontset=fandol]{ctex}
\RequirePackage{fontspec}
\setmainfont{cmunrm.otf}[BoldFont=cmunbx.otf,ItalicFont=cmunti.otf,BoldItalicFont=cmunbi.otf]
\setCJKfamilyfont{sourceheader}[ItalicFont=FandolKai-Regular.otf]{FandolKai-Regular.otf}
\RequirePackage{amsmath,amssymb,mathrsfs}
\RequirePackage{amsthm,etoolbox}
\RequirePackage{xcolor,geometry,setspace,fancyhdr,enumitem,needspace,xurl}
\definecolor{structurecolor}{HTML}{880E4F}
\definecolor{main}{HTML}{9C27B0}
\geometry{paperwidth=612bp,paperheight=792bp,left=20mm,right=20mm,top=95.04bp,bottom=20mm,headheight=12pt,headsep=25pt,footskip=30pt}
\setstretch{1.6}
\setlength{\parindent}{2em}
\setlength{\parskip}{0pt}
\setlist{nosep,leftmargin=2em}
\AtBeginDocument{\setlength{\abovedisplayskip}{4pt plus 1pt minus 1pt}\setlength{\belowdisplayskip}{4pt plus 1pt minus 1pt}\setlength{\abovedisplayshortskip}{2pt}\setlength{\belowdisplayshortskip}{3pt}}
\fancyhf{}
\fancyhead[L]{{\itshape\CJKfamily{sourceheader}\rightmark}}
\fancyhead[R]{{\itshape\CJKfamily{sourceheader}\leftmark}}
\fancyfoot[C]{\thepage}
\renewcommand{\headrulewidth}{0.4pt}
\renewcommand{\headrule}{{\color{structurecolor}\hrule height\headrulewidth width\headwidth\vskip-\headrulewidth}}
\pagestyle{fancy}
\fancypagestyle{cover}{\fancyhf{}\fancyfoot[C]{\thepage}\renewcommand{\headrulewidth}{0pt}\renewcommand{\headrule}{}}
\RequirePackage{hyperref}
\hypersetup{unicode=true,colorlinks=true,linkcolor=structurecolor,urlcolor=structurecolor,citecolor=structurecolor,linktoc=section,bookmarksnumbered=true,bookmarksopen=true,pdfborder={0 0 0},pdfstartview=FitH}
\setcounter{tocdepth}{1}
\renewcommand*\l@section{\@dottedtocline{1}{1.5em}{2.4em}}
\renewcommand{\thesection}{\thechapter.\arabic{section}}
\newcommand{\originalchapter}[1]{\par\addvspace{14pt}\Needspace{5\baselineskip}\refstepcounter{chapter}\setcounter{section}{0}\addcontentsline{toc}{chapter}{\protect\numberline{\thechapter}#1}\markboth{\thechapter\quad #1}{}\begingroup\centering\color{structurecolor}\bfseries\fontsize{14.4}{20}\selectfont\thechapter\quad #1\par\endgroup\nobreak\vspace{8pt}}
\newcommand{\originalsection}[1]{\par\addvspace{9pt}\Needspace{4\baselineskip}\refstepcounter{section}\addcontentsline{toc}{section}{\protect\hspace*{1.5em}\protect\numberline{\protect\textcolor{structurecolor}{\thesection}}#1}\markright{\thesection\quad #1}\noindent{\fontsize{12}{17}\selectfont\color{structurecolor}\thesection\quad}{\bfseries\fontsize{12}{17}\selectfont #1}\par\nobreak\vspace{4pt}}

% Semantic environments are registered once here, not re-designed in manuscripts.
\newtheoremstyle{wzstatement}{6pt}{5pt}
  {\normalfont\kaishu}{0pt}{\normalfont\bfseries\color{structurecolor}}
  {}{.7em}{\thmname{#1}\thmnumber{ #2}\thmnote{\normalfont\ (#3)}}
\newtheoremstyle{wzdefinition}{6pt}{5pt}
  {\normalfont\songti}{0pt}{\normalfont\bfseries\color{structurecolor}}
  {}{.7em}{\thmname{#1}\thmnumber{ #2}\thmnote{\normalfont\ (#3)}}
\newtheoremstyle{wzproblem}{7pt}{5pt}
  {\normalfont\songti}{2em}{\normalfont\bfseries\color{main}}
  {}{.7em}{\thmname{#1}\thmnumber{ #2}}
\newtheoremstyle{wzremark}{4pt}{4pt}
  {\normalfont\songti}{0pt}{\normalfont\bfseries}
  {}{.7em}{\thmname{#1}}
\theoremstyle{wzstatement}
\newtheorem{theorem}{定理}[chapter]
\newtheorem{lemma}[theorem]{引理}
\newtheorem{proposition}[theorem]{命题}
\newtheorem{corollary}[theorem]{推论}
\theoremstyle{wzdefinition}
\newtheorem{definition}[theorem]{定义}
\theoremstyle{wzproblem}
\newtheorem{example}{例题}[chapter]
\newtheorem{exercise}{习题}[chapter]
\theoremstyle{wzremark}
\newtheorem*{remark}{注}
\renewcommand{\proofname}{证明}
% Retain the native amsthm QED stack; only adjust the fixed typography.
\expandafter\patchcmd\csname\string\proof\endcsname{\normalfont \topsep6\p@\@plus6\p@\relax}
  {\normalfont\songti \topsep5\p@\@plus1\p@\relax}
  {}{\ClassError{elegantbook-original-adapter}{Cannot patch proof spacing}{Check the amsthm version.}}
\expandafter\patchcmd\csname\string\proof\endcsname{\itshape}{\normalfont\bfseries}
  {}{\ClassError{elegantbook-original-adapter}{Cannot patch proof heading}{Check the amsthm version.}}
\expandafter\patchcmd\csname\string\proof\endcsname{\ignorespaces}
  {\setlength{\parindent}{2em}\ignorespaces}
  {}{\ClassError{elegantbook-original-adapter}{Cannot patch proof paragraphs}{Check the amsthm version.}}
\AtBeginEnvironment{proof}{\par\Needspace{3\baselineskip}}
\renewcommand{\qedsymbol}{\ensuremath{\square}}
\newenvironment{solution}{\begin{proof}[解答]}{\end{proof}}
\newcommand{\wzcheckchapter}{\ifnum\value{chapter}<1
  \ClassError{elegantbook-original-adapter}{Numbered mathematics before chapter 1}{Move it after the first originalchapter.}\fi}

\newcommand{\wzregisterguard}[1]{\AtBeginEnvironment{#1}{\wzcheckchapter\par\Needspace{3\baselineskip}}}
\forcsvlist{\wzregisterguard}{theorem,lemma,proposition,corollary,definition,example,exercise}
% Front matter and bibliography are not numbered mathematical chapters.
\newcommand{\originalpreface}{\par\addvspace{10pt}\Needspace{5\baselineskip}
  \noindent{\bfseries\fontsize{12}{17}\selectfont 前言}\par\nobreak\vspace{4pt}}
\newcommand{\originalcontents}{\par\addvspace{14pt}
  \begin{center}{\color{structurecolor}\bfseries\fontsize{14.4}{20}\selectfont 目录}\end{center}
  \vspace{-10pt}\@starttoc{toc}}
\renewcommand{\bibname}{参考文献}
\renewenvironment{thebibliography}[1]{
  \par\addvspace{12pt}\Needspace{3\baselineskip}\phantomsection
  \addcontentsline{toc}{chapter}{\bibname}
  \markboth{\bibname}{}
  \noindent{\color{structurecolor}\bfseries\fontsize{14.4}{20}\selectfont\bibname}\par\nobreak\vspace{6pt}
  \list{\@biblabel{\@arabic\c@enumiv}}
    {\settowidth\labelwidth{\@biblabel{#1}}\leftmargin\labelwidth
     \advance\leftmargin\labelsep\usecounter{enumiv}
     \let\p@enumiv\@empty\renewcommand\theenumiv{\@arabic\c@enumiv}
     \setlength{\itemsep}{3pt}\setlength{\parsep}{0pt}}
  \sloppy\clubpenalty4000\widowpenalty4000\sfcode`\.=1000\relax
}{\def\@noitemerr{\@latex@warning{Empty bibliography}}\endlist}

\numberwithin{equation}{chapter}
\renewcommand{\theequation}{\thechapter.\arabic{equation}}
\allowdisplaybreaks[2]
\emergencystretch=2em
\clubpenalty=6000
\widowpenalty=6000
\raggedbottom
\endinput
\end{filecontents*}
% WZ-LOCKED-CLASS-END
\documentclass{elegantbook-original-adapter}
% ===== 元数据内容槽 =====
\newcommand{\BookTitle}{图论: 结构、染色与网络算法}
\newcommand{\BookAuthor}{MIT OCW 多课程资料; 中文重构与补证}
\newcommand{\BookDate}{18.315 Spring 2005 图论主干及指定补充}
\hypersetup{pdftitle={图论: 结构、染色与网络算法},pdfauthor={MIT OCW 原课程作者; 中文整理}}
\usepackage{tikz,booktabs,longtable}
\usetikzlibrary{arrows.meta,calc,positioning}
\newcommand{\R}{\mathbb R}
\newcommand{\N}{\mathbb N}
\newcommand{\Z}{\mathbb Z}
\newcommand{\E}{\mathbb E}
\newcommand{\Prob}{\mathbb P}
\newcommand{\rank}{\operatorname{rank}}
\newcommand{\vol}{\operatorname{vol}}
\newcommand{\dist}{\operatorname{dist}}
\newcommand{\supp}{\operatorname{supp}}
\newcommand{\eps}{\varepsilon}
\newcommand{\sourceof}[1]{\par\noindent\textit{来源定位: #1.}\par}
\tikzset{vertex/.style={circle,draw,fill=white,inner sep=1.5pt,minimum size=5mm},edge/.style={thick},arc/.style={-{Stealth[length=2mm]},thick}}
\begin{document}
\frontmatter
\begin{center}
{\color{structurecolor}\bfseries\fontsize{18}{25}\selectfont\BookTitle\par}
{\bfseries MIT OpenCourseWare 中文重构本\par}
\BookAuthor\par
\BookDate\par
\end{center}
\originalpreface
图用顶点和边记录对象之间的关系. 一条边的局部限制, 可以决定整个网络是否连通、是否存在匹配、需要多少种颜色, 以及信息和流量能否经过网络. 本书面向掌握集合、归纳法和基本线性代数的读者. 前十二章不预设图论知识; 矩阵树定理、电网络与谱方法需要实对称矩阵的特征值和行列式, 所用补充在相应位置交代. 概率方法只用有限样本空间、期望和条件概率.

主源为 Igor Pak 讲授、Amanda Redlich 记录的 \href{https://ocw.mit.edu/courses/18-315-combinatorial-theory-introduction-to-graph-theory-extremal-and-enumerative-combinatorics-spring-2005/pages/lecture-notes/}{18.315, Spring 2005} 中的图论主干. 原稿是手写课堂记录, 常把困难步骤交给作业或外部书籍; 本书按数学依赖重新组织, 在本地补齐核心证明. 基础和初等平面图来自 Lehman、Leighton、Meyer 的 6.042J Spring 2015 开放教材指定部分; 匹配与流算法补用 Santosh Vempala 的 18.433 Fall 2003; 树计数、电网络和有向 Euler 圈补用 Alexander Postnikov 讲授、Andrew Lin 记录的 18.212 Spring 2019; 极值与谱工具补用 Yufei Zhao 及该班学生的 18.217 Fall 2019 指定段落. 各来源重复定义与定理在正文合并.

这里的完整性指已声明的图论学习主线, 不指对上述整门组合课程逐讲全译. 连通性、树、Euler 与 Hamilton、匹配、顶点染色、边染色、平面图、多项式、搜索与最短路、最小生成树、网络流、矩阵树、电网络、谱方法、极值与 Ramsey 均形成从定义到证明和应用的连续单元. 逻辑、一般计数、一般概率、Young 表、分拆、纽结、加法组合和研究级图极限不属于本册. 四色定理、Kuratowski 判据及四连通平面图的 Hamilton 性只作为有明确边界的外部结果提及, 不冒称给出了其长篇证明. 书末来源对应表逐项区分采用、补证、选读和范围外内容.

中文正文为重新组织的数学改编, 不是原作者的逐字译稿. 例题和练习是经典数学问题的重新表述或本稿设置的模型, 未核实最初题源时不宣称原创. 图形均为针对本稿对象的原生 TikZ, 不复制外书图形、照片或受限题面. 第三方或权利不清的材料只保留链接. MIT 18.225 Fall 2023 的新版作者书稿因含独立版权及图片许可说明, 此次仅登记链接与核查元数据, 不打包其正文或插图.

中文改编依 \href{https://creativecommons.org/licenses/by-nc-sa/4.0/}{CC BY-NC-SA 4.0} 供非商业使用及相同方式共享. MIT OCW 默认条款见 \href{https://ocw.mit.edu/pages/privacy-and-terms-of-use/}{官方使用条款}; 6.042 原件明示 CC BY-NC-SA 3.0, 原件保留原许可, 改编依其第4(b)条采用同要素后续版本. 原作者与 MIT 均未为本项目背书. 逐文件页数、作者、年份、SHA-256 与版权核查记录见源码附带来源清单. 版本日期: 2026-10-08.
\clearpage
\originalcontents
\clearpage
\mainmatter
\originalchapter{图的语言与局部计数}\label{ch:basic}
\sourceof{6.042J 教材第11章; 18.315 图论部分的先修补全}
\originalsection{顶点、边与子图}
一张交通示意图可以变形而不改变道路连接关系. 图论首先保留这种关系, 再研究从关系本身能推出什么. 画在纸上的线段相交, 并不自动产生一个新顶点; 顶点必须在数学对象中指定.
\begin{definition}
有限简单无向图是 $G=(V,E)$, 其中 $V$ 为有限顶点集, $E$ 是 $V$ 的二元子集组成的集合. 边 $\{u,v\}$ 简记 $uv$, 其两个端点相邻. 简单图没有自环或平行边. 记 $n=|V|$, $m=|E|$, $N(v)=\{u:uv\in E\}$, $d(v)=|N(v)|$, 非空图的最小度与最大度分别为 $\delta(G)$、$\Delta(G)$. 空图约定最大度、顶点染色数及边染色数均为0, 最小度不定义.
\end{definition}
多重图允许两顶点间有若干条不同的边; 若还允许自环, 度数中每条自环计两次. 有向图的边为有序对, 称为弧; 顶点的入度和出度分别记 $d^-(v)$、$d^+(v)$. 除另行声明, 结构与染色章节均讨论有限简单无向图. 缩边、平面图对偶和流网络会显式允许平行边.

子图 $H=(W,F)$ 满足 $W\subseteq V$, $F\subseteq E$ 且各边端点属于 $W$. 诱导子图 $G[W]$ 保留 $W$ 内全部原边. 删除顶点 $v$ 时也删除与它关联的所有边; 删除边则不删端点. 补图 $\overline G$ 与 $G$ 顶点集相同, 两不同顶点在补图相邻当且仅当原图不相邻. 图同构是保持邻接关系的顶点双射, 因而编号与图形形状不是同构不变量.
\begin{example}
六个顶点组成一个六边形, 与两个互不相交的三角形, 是否同构?
\end{example}
\begin{solution}
两图都满足 $n=m=6$, 且每个顶点的度为2, 所以这些计数无法区分. 但前图只有一个连通分支, 后图有两个; 同构必须保持任意两顶点能否由路连接, 因而两图不同构. 度数列是有用的必要条件, 不是充分条件.
\end{solution}
完全图 $K_n$ 的任意两顶点相邻. 路图 $P_n$ 有 $n$ 个依次相连的顶点; 圈图 $C_n$ 在 $P_n$ 两端再加一条边, 其中 $n\ge3$. 完全二部图 $K_{a,b}$ 把顶点分为大小 $a,b$ 的两部分, 只在不同部分之间连边, 且所有这种边都存在.
\originalsection{握手引理与度数列}
\begin{theorem}[握手引理]\label{thm:handshake}
有限无向多重图满足 $\sum_{v\in V}d(v)=2m$. 有向图满足 $\sum_v d^+(v)=\sum_v d^-(v)=m$.
\end{theorem}
\begin{proof}
按关联的端点计数每条无向边, 非自环贡献两个不同端点, 自环在同一端点贡献两次. 有向弧恰有一个起点和一个终点, 对两种度分别求和均计每条弧一次.
\end{proof}
\begin{corollary}
任意有限无向图中, 奇度顶点的个数为偶数.
\end{corollary}
\begin{proof}
在度数总和中删去所有偶数项, 剩下若干奇数之和仍为偶数; 奇数项数只能为偶数.
\end{proof}
平均度是 $\overline d=2m/n$ ($n>0$). 它不意味着每个顶点都有相近的度, 但能保证某个顶点度不超过它、某个顶点度不小于它. 平面图的染色和极值图论都反复使用这个转换.
\begin{theorem}[Havel--Hakimi 归约]\label{thm:hh}
设 $d_1\ge d_2\ge\cdots\ge d_n\ge0$ 为整数, $d_1=k\le n-1$. 它是简单图的度数列, 当且仅当从 $d_2,\ldots,d_{k+1}$ 各减1, 保留其余项并重新排序所得序列是简单图的度数列. 归约出现负数时不可实现.
\end{theorem}
\begin{proof}
归约序列若可实现, 加一个新顶点并连接到被减1的 $k$ 个顶点即可. 反向需要说明为什么能指定最高度顶点的邻居. 在一张实现原序列的图中, 设最高度顶点为 $v$. 令 $T$ 为预定排序中的前 $k$ 个其他顶点. 若 $v$ 尚未邻接全部 $T$, 取 $x\in T\setminus N(v)$、$y\in N(v)\setminus T$, 于是 $d(x)\ge d(y)$. 我们把邻点 $y$ 换成 $x$.

具体说, 若找不到 $z\notin\{v,x,y\}$ 使 $xz\in E$ 而 $yz\notin E$, 则 $x$ 除可能的 $y$ 外的邻居都是 $y$ 的邻居; $y$ 还多出邻居 $v$, 导致 $d(y)>d(x)$, 矛盾. 取这样的 $z$, 删除 $vy,xz$, 加入 $vx,yz$. 新边原先不存在, 四个端点互异, 度数全不变. 每次交换使 $|N(v)\cap T|$ 增加1, 有限次后 $v$ 邻接指定的前 $k$ 个顶点. 删除 $v$ 得到归约序列.
\end{proof}
\begin{example}
判断 $(3,3,2,2,2)$ 与 $(3,3,3,1)$ 是否可实现.
\end{example}
\begin{solution}
第一列归约为 $(2,2,1,1)$, 再归约为 $(1,1,0)$, 最后为 $(0,0)$, 所以可实现. 第二列归约为 $(2,2,0)$, 再归约出现负数, 所以不可实现. 第二列的度数和虽为偶数, 但三个度3顶点都必须连接度1顶点, 已经矛盾. 握手条件本身不够.
\end{solution}
\originalsection{行走、路、圈与连通}
\begin{definition}
行走是相邻顶点组成的序列 $v_0,v_1,\ldots,v_\ell$, 长度为 $\ell$. 迹不重复边, 路不重复顶点. 闭行走满足 $v_0=v_\ell$; 圈除首尾重合外不重复顶点. 无向简单图的圈长度至少3. 两顶点之间有路时称它们连通; 极大的相互连通顶点集称为连通分支.
\end{definition}
从行走中删去两次出现同一顶点之间的闭段, 可得到同端点的路. 因而定义连通性时用行走或路等价. “有闭行走”与“有圈”则不同: 一条边来回走形成闭行走, 却没有圈. 若闭行走长度为奇数, 它一定含奇圈: 在最短奇闭行走中, 若有内部重复顶点, 分成的两个闭段之一仍为奇长, 与最短性矛盾.
\begin{theorem}\label{thm:bipartite}
无向图为二部图, 当且仅当它不含奇圈.
\end{theorem}
\begin{proof}
二部图的行走每一步都换到另一部分, 所以圈长度为偶数. 反向对每个分支选根 $r$, 按到根的最短路长度的奇偶分成两部分. 若某条边 $uv$ 的端点同属一部分, 从根到 $u$ 的最短路、边 $uv$ 和从 $v$ 回根的最短路组成奇闭行走, 因而含奇圈, 矛盾. 所有边都跨两部分.
\end{proof}
对连通图, 距离 $\dist(u,v)$ 是两顶点间最短路长度, 直径是所有距离的最大值. 圈图 $C_n$ 的直径为 $\lfloor n/2\rfloor$, 完全图 $K_n$ ($n\ge2$) 的直径为1. 非连通图可将不同分支之间的距离记为 $+\infty$; 讨论有限直径时必须先要求连通.
\originalsection{习题与解答}
\begin{exercise}\label{ex:basic1}
证明非空简单图中至少有两个顶点度数相同, 其中顶点数 $n\ge2$.
\end{exercise}
\begin{exercise}\label{ex:basic2}
证明 $n$ 个顶点、$m$ 条边的图有独立集大小至少 $n/(\Delta+1)$, 其中独立集指内部没有边的顶点集.
\end{exercise}
习题\ref{ex:basic1}的解答.
\begin{solution}
度数在 $0,1,\ldots,n-1$ 中. 度0和度 $n-1$ 不可能同时出现, 所以至多有 $n-1$ 个可取值, 对 $n$ 个顶点应用抽屉原理.
\end{solution}
习题\ref{ex:basic2}的解答.
\begin{solution}
不断取一个尚未删除的顶点放入独立集, 并删除它和所有邻居. 每次至多删除 $\Delta+1$ 个顶点; 选取的顶点彼此不邻接. 删除完全部顶点至少需 $n/(\Delta+1)$ 次. 整数意义下可向上取整.
\end{solution}

\originalchapter{树与生成树}\label{ch:trees}
\sourceof{6.042J 第11.10节; 18.212 L22图论部分、L23}
\originalsection{树的等价刻画}
树是连通而没有圈的图. 它没有多余的连接, 又能连接全部顶点, 所以在网络建设、递归结构和计数中都自然出现. 一个没有圈的图称为森林, 每个分支都是树; 单顶点也算树.
\begin{lemma}\label{lem:leaf}
含至少两个顶点的有限树至少有两个度1顶点, 称为叶子.
\end{lemma}
\begin{proof}
取最长路 $v_0\cdots v_\ell$. 端点 $v_0$ 若有路外邻居就能延长路; 若有除 $v_1$ 之外的路内邻居就产生圈. 因而它只有一个邻居. 另一端点同理.
\end{proof}
\begin{theorem}\label{thm:tree-equivalent}
对 $n\ge1$ 的简单图 $G$, 以下条件等价: (1) $G$ 是树; (2) 任意两顶点之间恰有一条路; (3) $G$ 连通且 $m=n-1$; (4) $G$ 无圈且 $m=n-1$; (5) $G$ 连通, 删除任意边都不连通.
\end{theorem}
\begin{proof}
树中两顶点有路. 两条不同的路沿首次分开再会合的区段产生圈, 所以路唯一. 唯一路的任一边删除后, 其两端不能再连通, 得到(5). 反向, 若连通图有圈, 删除圈上一边仍可沿圈另一侧通行, 违反(5).

由叶子引理, 对树删去一片叶子得到小一阶的树, 从单顶点归纳得到 $m=n-1$. 任意连通图都可不断删除圈上的边而保持连通, 直到得到一棵包含全部顶点的树, 称为生成树. 它有 $n-1$ 条边. 若原图本来只有 $n-1$ 条边, 就无可删除的边, 原图本身为树, 得到(3)的反向. 森林若有 $c$ 个分支, 在各分支分别计数得到 $m=n-c$; (4)使 $c=1$. 这同时完成全部等价关系.
\end{proof}
树加一条新边, 恰好产生一个圈: 它由新边与树中连接其端点的唯一路组成. 树删一条边恰好分成两个分支: 端点已不连通, 而每个顶点沿原来的路总能到达至少一个端点. 后面的最小生成树算法将围绕这两个操作构造交换论证.
\begin{example}
一棵树中度1、2、3的顶点个数分别为 $a,b,c$, 且没有其他度数. 求 $a$ 与 $c$ 的关系.
\end{example}
\begin{solution}
若树有至少两个顶点, 握手引理给出 $a+2b+3c=2(a+b+c-1)$, 所以 $a=c+2$. 度2顶点只把边细分, 不影响分枝与叶子之间的数量关系. 单顶点度0, 已不满足题面所列度数种类.
\end{solution}
\originalsection{根树与生成森林}
指定根后, 每个非根顶点在通向根的唯一路上有一个父顶点, 其余相邻顶点为子顶点. 顶点的深度是到根的距离. 从根向外逐层访问得到的父边必组成生成树, 但不同访问策略可能得到不同树.

在一般图中, 极大的无圈边集必在每个原分支内连通: 若一个原分支被分为两块, 可加入跨两块的边而不产生圈, 与极大性矛盾. 因而它是生成森林, 有 $n-c(G)$ 条边. 注意“极大”是不能再包含更多边, “最大”是边数达到最大; 这里因为森林的结构, 二者最终边数相同, 不能在其他优化问题中随便混用.
\begin{center}
\begin{tikzpicture}[x=1.25cm,y=1cm]
\node[vertex](r)at(0,0){$r$};\node[vertex](a)at(-1,-1){$a$};\node[vertex](b)at(1,-1){$b$};
\node[vertex](c)at(-1.6,-2){$c$};\node[vertex](d)at(-.4,-2){$d$};\node[vertex](e)at(1,-2){$e$};
\foreach \u/\v in {r/a,r/b,a/c,a/d,b/e}\draw[edge](\u)--(\v);
\draw[dashed,bend right=35](c)to(e);
\node at(3,-.6){实线构成生成树};\node at(3,-1.2){加虚线产生唯一圈};
\end{tikzpicture}
\end{center}
图中虚线若加入, 圈为 $c,a,r,b,e,c$. 从这条圈中删去任一原树边, 就得到另一棵生成树. 没在圈上的边不能代替它被删, 因为那样仍有圈且另有部分被断开.
\originalsection{Prüfer 编码与 Cayley 公式}
\begin{theorem}[Cayley 公式]\label{thm:cayley}
顶点集为 $[n]=\{1,\ldots,n\}$ 的树共有 $n^{n-2}$ 棵, 其中 $n\ge2$.
\end{theorem}
\begin{proof}
给每棵标号树作 Prüfer 编码: 每次删除当前编号最小的叶子, 记录其邻居编号, 直到只剩两个顶点. 得到长度 $n-2$ 的序列 $(a_1,\ldots,a_{n-2})$, 每项在 $[n]$ 中.

反向给定任意这种序列. 在尚未删除的标号中, 取未在当前剩余序列出现的最小标号 $b_i$, 加边 $b_i a_i$, 删去 $b_i$ 并删掉首项 $a_i$. 该标号存在, 因为当前顶点数比序列长度多2. 序列中出现的标号不会提前被删除, 所以 $a_i$ 是一个仍在的、不同于 $b_i$ 的顶点. 最后连接剩下的两顶点.

构造出的图连通: 每个删除顶点都连到一个以后才删除的顶点, 顺着删除次序最终到达最后一条边. 它有 $n-1$ 条边, 因而是树. 在当前构造树中, 一个顶点的度数是它在当前序列中的出现次数加1: 每次作为 $a_i$ 接收一条边, 最后作为 $b_i$ 或最后两点之一又有一条边. 所以当前叶子恰是序列中不出现的标号, 两个过程每一步选出的最小叶子一致. 编码与解码互逆, 树数等于序列数 $n^{n-2}$.
\end{proof}
\begin{corollary}
若正整数 $d_1,\ldots,d_n$ 满足 $\sum d_i=2n-2$, 则具有该指定度数的标号树共有
\[
\frac{(n-2)!}{\prod_{i=1}^n(d_i-1)!}.
\]
\end{corollary}
\begin{proof}
由编码证明, 标号 $i$ 在 Prüfer 序列中恰出现 $d_i-1$ 次. 这些次数之和为 $n-2$, 按多重排列计数即得. 条件也使这样的序列确实存在.
\end{proof}
\begin{example}
解码序列 $(4,4,2,4)$, 并求树的度数.
\end{example}
\begin{solution}
这里 $n=6$. 首先缺失的最小标号为1, 加边 $14$; 剩余序列 $(4,2,4)$ 中缺失的最小尚存标号为3, 加边 $34$. 然后对 $(2,4)$ 取5, 加边 $52$; 对 $(4)$ 取2, 加边 $24$; 最后连接4、6. 度数为 $(1,2,1,4,1,1)$, 与每个标号出现次数加1一致.
\end{solution}
\originalsection{习题与解答}
\begin{exercise}\label{ex:tree1}
证明任意树都是二部图, 并说明它何时存在完美匹配.
\end{exercise}
\begin{exercise}\label{ex:tree2}
顶点为 $[n]$, 指定顶点1的度恰为 $k$, 其中 $1\le k\le n-1$. 求树的数量.
\end{exercise}
习题\ref{ex:tree1}的解答.
\begin{solution}
树没有圈, 特别没有奇圈, 由定理\ref{thm:bipartite}为二部图. 完美匹配可以从叶子判断: 叶子 $u$ 的唯一边 $uv$ 必须被使用, 删去 $u,v$ 后, 原树有完美匹配当且仅当剩余森林有完美匹配. 反复执行, 最终为空图则存在, 出现孤立顶点则不存在. 偶数阶仅是必要条件, 例如四叶星有五点, 而五叶星有六点也无完美匹配. 匹配的正式理论见第\ref{ch:matching}章.
\end{solution}
习题\ref{ex:tree2}的解答.
\begin{solution}
在长度 $n-2$ 的序列中选 $k-1$ 个位置填1, 其余位置任填另外 $n-1$ 个标号. 数量为 $\binom{n-2}{k-1}(n-1)^{n-k-1}$, 边界 $k=n-1$ 时只有全1序列, 对应以1为中心的星.
\end{solution}

\originalchapter{割、连通度与块}\label{ch:connectivity}
\sourceof{18.315 L23; 块分解与排序细节为编者补充}
\originalsection{边割与顶点割}
连通只说明还能通行, 没有衡量网络对故障的承受能力. 删除一条边或一个顶点可能产生完全不同的结果. 顶点故障连带删除所有关联边, 因而需要区分两种连通度.
\begin{definition}
对 $\varnothing\ne S\subsetneq V$, 边割 $\partial S$ 是一端在 $S$、另一端在补集的边集. 连通图中删去后不再连通的单条边称为桥. 若删去顶点 $v$ 增加分支数, 则 $v$ 是割点. 连通非完全图的顶点连通度 $\kappa(G)$ 是使图不连通所需删除的最少顶点数; 完全图约定 $\kappa(K_n)=n-1$. 边连通度 $\lambda(G)=\min|\partial S|$ ($n\ge2$). 非连通图的两种连通度均记0.
\end{definition}
\begin{proposition}
一条边是桥当且仅当它不在任何圈上. 对 $n\ge2$ 的连通图, $\kappa(G)\le\lambda(G)\le\delta(G)$.
\end{proposition}
\begin{proof}
若边在圈上, 圈的另一段提供替代通路, 故不是桥. 若非桥, 删边后其两端仍由路相连, 加回该边就产生圈. 单顶点割的边数为它的度, 所以 $\lambda\le\delta$.

设最小边割分开非空 $A,B$, 其大小为 $k$. 如果存在 $a\in A,b\in B$ 不相邻, 对每条跨割边选择一个端点删除, 避免选择 $a,b$: 含 $a$ 的边删另一端, 含 $b$ 的边亦然, 其余任选. 至多删 $k$ 个顶点后 $a,b$ 分离. 若两侧完全相连, 则 $k=|A||B|\ge |A|+|B|-1=n-1\ge\kappa$. 得到所需不等式.
\end{proof}
一条桥的两端不一定都是割点: 单边图的任何顶点删除后只剩单点, 分支数没有增加. 因此把桥与割点混为一谈会在小图上立即出错.
\originalsection{不相交路与 Menger 定理}
\begin{theorem}[Menger 的边版本]\label{thm:mengeredge}
在有限无向图中, 两不同顶点 $s,t$ 之间两两边不相交路的最大条数, 等于使 $s,t$ 分离的最小边割大小.
\end{theorem}
\begin{proof}
每条 $s$--$t$ 路都必须穿过边割, 边不相交路不能共用该割中的边, 所以路数不超过割大小. 对反向, 把每条无向边换成两个反向弧, 容量各为1. 第\ref{ch:flow}章用独立的增广算法证明整数最大流最小割定理; 一个无向割的有向前向容量恰等于其边数. 取整数最大流, 若原边两个方向均有流量, 同时减去相同流量, 不改变任何顶点净流量. 剩余每条原边最多使用一个方向. 反复沿正流量从 $s$ 找到 $t$, 删去一单位路流; 若出现圈先消去圈流. 于是最大流分解为边不相交路, 数目等于最小割. 流定理的证明不依赖本章, 所以这里没有循环引用.
\end{proof}
\begin{theorem}[Menger 的顶点版本]\label{thm:mengervertex}
设 $s,t$ 为不相邻的不同顶点. 内部顶点两两不相交的 $s$--$t$ 路的最大条数, 等于可从 $V\setminus\{s,t\}$ 删除以分离 $s,t$ 的最小顶点数.
\end{theorem}
\begin{proof}
每条路必须经过所删顶点之一, 因而得到上界. 为证明下界, 把每个非端点 $v$ 拆为 $v_{\rm in},v_{\rm out}$, 加容量1的弧 $v_{\rm in}\to v_{\rm out}$. 原边按可行方向变成从出端到入端的弧, 容量取 $M=n+1$, 端点 $s,t$ 不拆. 不相邻条件使删除全部其他顶点能分离它们, 所以存在容量至多 $n-2$ 的割. 最小割不可能含容量 $M$ 的弧, 只能由某些容量1的顶点弧组成. 这些顶点在原图构成分离集, 反之每个分离集给出不大于其大小的割.

整数最大流分解为路后, 每条容量1的顶点弧最多被一路使用, 投影回原图就是内部顶点不相交路. 删除投影中的可能重复段不增加共用顶点. 因而最大路数等于最小顶点割.
\end{proof}
“不相邻”不能删去: 若 $st$ 本来是一条边, 删除其他顶点永远无法分开 $s,t$, 但内部不相交路数仍有限. 对 $n\ge k+1$, $1\le k\le n-1$, 全局 $k$ 连通性等价于任意两点之间有 $k$ 条内部顶点不相交路. 非相邻端点的正向由 Menger 立即得到. 对相邻的 $s,t$, 若集合 $S\subseteq V\setminus\{s,t\}$ 在 $G-st$ 中分离它们且 $|S|<k-1$, 令 $A$ 为 $(G-st)-S$ 中含 $s$ 的分支. 若 $A\setminus\{s\}\ne\varnothing$, 删除 $S\cup\{s\}$ 后该集合仍与 $t$ 分离, 得大小小于 $k$ 的顶点割. 若 $A=\{s\}$, 则 $d_G(s)\le|S|+1<k$, 也与 $k$ 连通图最小度至少 $k$ 矛盾. 故删边图的最小分离集至少有 $k-1$ 点, 用非相邻版本再加单边路 $st$. 反向, 删除少于 $k$ 个顶点时, 对任意两个未删点, 至多截断相应数量的不相交路, 至少一路保留, 所以余图连通.

\originalsection{块树}
Brooks 染色定理需要在图中选两顶点一起删除, 仍保留连通性. 仅知道“无割点”不直接保证任意两点都能删除, 因而需要描述有割点图怎样拼接.
\begin{definition}
连通图的块包括极大二连通子图和每条桥组成的二顶点子图; 单顶点图单独作为一个块. 二连通图有至少三个顶点, 连通且无割点. 块的内部顶点指在该块中、却不是原图割点的顶点.
\end{definition}
\begin{lemma}\label{lem:blocktree}
连通图的不同块至多共有一个顶点; 有共有顶点时它是割点. 以块和割点为两类节点, 当割点属于某块就连边, 得到的关联图为树. 若原图有割点, 该树至少有两个叶块; 每个叶块恰含一个割点, 删除其任一内部顶点仍使原图连通.
\end{lemma}
\begin{proof}
先说明二连通子图的粘合性质. 若两个二连通子图至少共有两点, 删任一点后仍至少有一个共有点, 两子图的剩余部分都连通, 所以并图无割点且连通. 它们若为极大块, 必属同一块. 桥不能与二连通子图共用两端, 否则桥在圈上. 故不同块至多共有一点.

若两块共点 $v$, 却有一条避免 $v$ 的路连接两块的其他顶点, 可取路径的首次、末次接触点, 并把路径与两块合并. 删除任一点后, 经过 $v$ 的连接与避开 $v$ 的路径至少保留一种, 各块内部也保持可连接; 若块是桥则其端点可沿这条新圈连接. 合并得到更大的无割点子图, 违反极大性. 所以共有点必为割点.

每条边属于某块: 非桥边在圈上, 从圈扩张到极大二连通子图即可. 原图路径可沿所属块转移, 关联图于是连通. 若关联图有圈, 其上每个块有两个不同的连接割点. 圈上块的并图删去任一点后, 各二连通块的剩余部分仍连通, 桥块至多失去一个端点. 圈上的不同连接点至多失去一个, 其余连接仍沿圈形成通路, 因而并图无割点, 会合成更大块, 矛盾. 因而关联图为树. 删去割点所得各分支都含它的邻点; 这些邻边不可能属于同一二连通块, 所以割点节点至少关联两个块, 所以叶节点只能是块. 有割点时树有至少两叶, 叶块只有一个连接割点. 删除内部顶点, 二连通叶块的余部仍连通; 若它是桥, 删除内部端点只去掉叶子. 其他块均不受影响.
\end{proof}
\begin{example}
两个三角形只共有一个顶点 $v$. 求它的块、割点、$\kappa$ 与 $\lambda$.
\end{example}
\begin{solution}
两个三角形为两个块, $v$ 为唯一割点, 块树是“块--$v$--块”三节点路. 因而 $\kappa=1$. 每条边在三角形上, 无桥, 故 $\lambda\ge2$; 顶点 $v$ 之外各点度2, 故 $\lambda=2$. 边冗余不能替代顶点冗余.
\end{solution}
\originalsection{习题与解答}
\begin{exercise}\label{ex:conn1}
求 $K_{a,b}$ 的顶点连通度和边连通度, 其中 $1\le a\le b$ 且 $a+b\ge3$.
\end{exercise}
\begin{exercise}\label{ex:conn2}
证明二连通图任意两条不同边都在同一个圈上.
\end{exercise}
习题\ref{ex:conn1}的解答.
\begin{solution}
删去小侧全部 $a$ 点后, 大侧有至少两点而不连通, 所以 $\kappa\le a$. 删去少于 $a$ 点后两侧仍非空, 完全二部图仍连通, 所以 $\kappa=a$. 又 $\kappa\le\lambda\le\delta=a$, 得 $\lambda=a$.
\end{solution}
习题\ref{ex:conn2}的解答.
\begin{solution}
把两条边分别细分, 加新顶点 $x,y$. 细分保持二连通性: 删除原顶点时, 被细分边的剩余段仍附着在另一端; 删除新顶点相当于删原边, 而二连通图无桥. 由 Menger 得两条内部不相交的 $x$--$y$ 路, 合起来构成经过 $x,y$ 的圈. 两新顶点度2, 圈必使用其两条关联边, 恢复原边后得到所求圈.
\end{solution}

\originalchapter{Euler 迹与 Hamilton 圈}\label{ch:eulerhamilton}
\sourceof{18.315 L19--22; Euler 先修补全及证明细节为编者补充}
\originalsection{一次走完所有边}
Euler 问题要求每条边恰走一次, Hamilton 问题要求每个顶点恰访问一次. 两者看似相近, 但控制对象不同. Euler 性由度数与连通性决定, Hamilton 性通常不能由一个同样简单的局部条件决定.
\begin{theorem}[无向 Euler 判据]\label{thm:euler}
有限无向多重图若至少有一条边, 存在使用全部边的闭迹当且仅当所有正度顶点处于同一分支且所有顶点度数为偶数. 存在端点不同的 Euler 迹当且仅当正度顶点处于同一分支且恰有两个奇度顶点; 这两个顶点就是端点.
\end{theorem}
\begin{proof}
闭迹每次进入一个顶点后都要离开, 在该点的关联边成对使用, 所以度数为偶数. 开迹只有两个端点各多一次进入或离开. 一条迹经过的所有边当然位于同一分支.

反向先考虑偶度情形. 从任一正度顶点出发, 不断沿未用边行走, 直到无边可走. 若停在不同于起点的顶点, 已用关联边数为奇数, 而原度为偶数, 仍应有未用边, 矛盾. 因而得到闭迹. 如果还有边没用, 正度部分的连通性保证闭迹上有顶点关联未用边: 可沿原图从未用边端点走向已有闭迹, 取首次接触点. 剩余未用边子图仍为偶度, 从该点再作闭迹, 将新闭迹插入旧迹. 每次用掉更多边, 有限次后用完全部边. 这就是 Hierholzer 构造. 恰有两奇点时, 在它们之间添加一条辅助边, 得到偶度图; 对闭迹删去辅助边并从该处切开, 得到所求开迹.
\end{proof}
孤立顶点不影响“走完所有边”. 若要求同时访问全部顶点, 就必须单独排除孤立点. 无边图则可把单顶点处的零长闭迹当作退化 Euler 迹, 与上面有边的判据区分.
\begin{theorem}[有向 Euler 判据]\label{thm:directedeuler}
有限有向多重图至少有一条弧, 且所有非零度顶点在忽略方向后连通. 若每点满足 $d^+=d^-$, 则有有向 Euler 闭迹. 若两点 $s,t$ 分别满足 $d^+(s)=d^-(s)+1$, $d^-(t)=d^+(t)+1$, 其余点平衡, 则有从 $s$ 到 $t$ 的有向 Euler 迹. 这些条件也都必要.
\end{theorem}
\begin{proof}
必要性由沿迹的进出计数得到. 平衡时沿未用出弧行走, 除起点外任何停点都有“已用入弧比出弧多1”, 与总入出度相等矛盾; 所以先得到闭迹. 剩余弧仍平衡. 若未用部分与已有闭迹相邻, 连接弧可以是入弧或出弧; 闭迹顶点的剩余入出度相等, 有任一未用入弧也必有未用出弧. 因而仍可在该顶点开始拼接新的有向闭迹. 弱连通性保证能找到这样的接触点. 非平衡情形添加 $t\to s$, 再切开闭迹.
\end{proof}
\originalsection{访问所有顶点}
\begin{definition}
经过图中所有顶点恰一次的圈称为 Hamilton 圈; 经过所有顶点恰一次的路称为 Hamilton 路. 只有至少三个顶点的简单图可能有 Hamilton 圈.
\end{definition}
\begin{proposition}\label{prop:ham-cut}
若 $G$ 有 Hamilton 圈, 则对任何非空 $S\subseteq V$, 图 $G-S$ 的分支数至多 $|S|$.
\end{proposition}
\begin{proof}
从 Hamilton 圈删除 $S$, 剩余部分至多为 $|S|$ 段路, 每段在 $G-S$ 中仍连通. 添加原图其余边只能合并分支, 不能增加分支. 若删去所有顶点, 分支数为0, 仍满足结论.
\end{proof}
\begin{example}
一个连通图有两个三角形只在一顶点相交. 它有 Euler 闭迹吗? 有 Hamilton 圈吗?
\end{example}
\begin{solution}
共有顶点度4, 其余顶点度2, 所以有 Euler 闭迹, 可顺次绕两个三角形. 但删共有顶点后有两个分支, 大于所删顶点数1, 故没有 Hamilton 圈. 在 Euler 迹中重复访问顶点是允许的.
\end{solution}
\begin{lemma}[闭包引理]\label{lem:closure}
简单图 $G$ 有 $n\ge3$ 个顶点. 若非相邻点 $u,v$ 满足 $d(u)+d(v)\ge n$, 则 $G$ 有 Hamilton 圈当且仅当 $G+uv$ 有 Hamilton 圈.
\end{lemma}
\begin{proof}
正向由加边立即成立. 反向若新增边未被 Hamilton 圈使用, 原圈就在 $G$ 中. 若被使用, 去掉它得 Hamilton 路 $x_1=u,x_2,\ldots,x_n=v$. 在索引集 $\{1,\ldots,n-1\}$ 中令
\[
A=\{i:ux_{i+1}\in E\},\qquad B=\{i:vx_i\in E\}.
\]
因为 $u,v$ 不相邻, 所有邻点都在路的内部, 有 $|A|=d(u)$、$|B|=d(v)$. 索引集只有 $n-1$ 项, $|A|+|B|\ge n$ 使两集合相交. 对相交索引 $i$, 圈 $u,x_2,\ldots,x_i,v,x_{n-1},\ldots,x_{i+1},u$ 使用全部顶点且所有边都在原图中.
\end{proof}
\begin{theorem}[Ore 与 Dirac]\label{thm:dirac}
若简单图 $G$ 有 $n\ge3$ 个顶点, 且每对非相邻顶点都满足 $d(u)+d(v)\ge n$, 则 $G$ 有 Hamilton 圈. 特别地, $\delta(G)\ge n/2$ 足够保证 Hamilton 性.
\end{theorem}
\begin{proof}
不断添加缺失边. 已有顶点的度只能增加, 因而闭包引理在每一步都可用. 最后得到有 Hamilton 圈的完全图 $K_n$, 再逐步反推原图也有. 最小度条件保证每对非相邻点的度数和至少为 $n$.
\end{proof}
这两个条件是充分条件, 并非必要条件. 圈图 $C_n$ 在 $n\ge5$ 时最小度2小于 $n/2$, 却本身就是 Hamilton 圈. 另一方面, 把两个大小尽可能均衡的团用一个割点粘合, 可以得到最小度接近 $n/2$ 而无 Hamilton 圈的图, 说明度数门槛的量级有实际意义.
\originalsection{平面图中的 Grinberg 必要条件}
\begin{theorem}[Grinberg 条件]\label{thm:grinberg}
设有限连通平面嵌入图有 Hamilton 圈 $C$. 记圈内、圈外边界长度为 $k$ 的面数分别为 $f_k^+,f_k^-$. 面的边界按行走计重数. 则
\[
\sum_{k\ge3}(k-2)(f_k^+-f_k^-)=0
\]
对简单二连通嵌入成立.
\end{theorem}
\begin{proof}
Hamilton 圈包含全部 $n$ 个顶点. 设内部弦数为 $e_+$, 内部面数为 $f_+$. 把圆盘外部当作一个面, 第\ref{ch:planar}章的 Euler 公式给 $n-(n+e_+)+(f_++1)=2$, 即 $f_+=e_++1$. 内部面边界长度总和为 $n+2e_+$: 圈边计一次, 弦计两次. 因而 $\sum(k-2)f_k^+=n+2e_+-2(e_++1)=n-2$. 外部区域在球面上也是圆盘, 完全同理得到 $\sum(k-2)f_k^-=n-2$. 相减即得.
\end{proof}
公式只能排除某些 Hamilton 圈, 满足它不意味着圈存在. 若图存在桥或重复面边界, 必须先检查 Hamilton 假设已排除这些障碍, 不能只凭一张画图机械代入.
\originalsection{奇度图的 Hamilton 圈奇偶性}
\begin{theorem}[奇度图的奇偶性]\label{thm:smith}
有限简单图中若每个顶点度数为奇数, 则每条固定边属于偶数个 Hamilton 圈. 特别地, 三正则图若有 Hamilton 圈, 就至少有三个不同的 Hamilton 圈.
\end{theorem}
\begin{proof}
若 $n\le2$, 没有 Hamilton 圈, 结论成立. 以下设 $n\ge3$, 固定边 $uv$, 此时 Hamilton 路的末点前驱与 $u$ 不同. 考虑所有以有序前缀 $(u,v)$ 开始的 Hamilton 路, 每条写作 $P=(u,v,x_3,\ldots,x_n)$. 构造辅助图, 它们为顶点. 若末点 $x_n$ 邻接 $x_i$ 且 $2\le i\le n-2$, 将路尾旋转为
\[
(u,v,\ldots,x_i,x_n,x_{n-1},\ldots,x_{i+1}).
\]
这是另一条保留前缀的 Hamilton 路, 旋转可逆; 不同的 $i$ 产生不同的新末点, 所以辅助图的度数就是可旋转索引数. 末点与前驱必相邻, 与 $u$ 的邻接不能拿来旋转, 故其辅助度为 $d_G(x_n)-1-\mathbf1_{ux_n\in E}$. 因原度为奇数, 辅助度为奇数当且仅当 $ux_n\in E$.

辅助图握手引理说明奇度顶点数为偶数. 它们正好对应把末点连回 $u$ 所得、包含 $uv$ 的 Hamilton 圈. 固定前缀确定了圈的走向, 每个无向圈恰对应一条这样的路, 所以圈数为偶数.

三正则图的一个 Hamilton 圈包含边 $e$, 因奇偶性还有另一个包含 $e$ 的圈. 两圈不同, 原圈中有边 $f$ 不在第二圈中. 再对 $f$ 用奇偶性, 得到第三个不同的圈.
\end{proof}
\originalsection{置换图的构造}
18.315 的后续例子把 Hamilton 性放在 Cayley 图中讨论. 这里保留一个能够从头构造的实例. 以 $S_n$ 的置换为顶点, 两置换若交换相邻位置即可互相得到就连边. 对 $n\ge3$, 这张图有 Hamilton 圈.

归纳列举全部置换: 已有 $(n-1)$ 阶置换的相邻交换列表时, 对第一项把 $n$ 从最右位置依次向左移动, 对第二项从最左位置向右移动, 如此交替. 一个块内相邻项只交换 $n$ 与其邻居; 两块之间 $n$ 在同一端, 其余项按旧列表交换相邻位置. 每个置换恰好属于一个块. 当 $n\ge3$ 时 $(n-1)!$ 为偶数, 最后一块与第一块中的 $n$ 都在最右端, 首尾其余置换可由一次相邻交换连接. 归纳基点 $n=3$ 的列表为 $123,132,312,321,231,213$, 首尾也相邻. 这是 Hamilton 圈的实际构造, 不依赖仅有连通性的判断.
\originalsection{习题与解答}
\begin{exercise}\label{ex:euler1}
给 $K_{a,b}$ 判断 Euler 闭迹与 Hamilton 圈的存在条件, $a,b\ge1$.
\end{exercise}
\begin{exercise}\label{ex:euler2}
证明若简单图 $G$ 有 $n\ge2$ 个顶点且 $\delta(G)\ge(n-1)/2$, 则有 Hamilton 路.
\end{exercise}
习题\ref{ex:euler1}的解答.
\begin{solution}
图连通, 大小为 $a$ 的一侧顶点度 $b$, 另一侧顶点度 $a$, 因而 Euler 闭迹当且仅当 $a,b$ 都为偶数. Hamilton 圈在两侧交替, 必须 $a=b$, 且简单圈要求 $a=b\ge2$; 此时依次交替使用各侧顶点即可构造.
\end{solution}
习题\ref{ex:euler2}的解答.
\begin{solution}
加一个与所有原顶点相邻的新顶点. 新图有 $n+1$ 个顶点, 原顶点度至少 $(n+1)/2$, 新点度 $n\ge(n+1)/2$. 用 Dirac 得 Hamilton 圈, 再删新点及其两条圈边, 留下原图的 Hamilton 路.
\end{solution}

\originalchapter{匹配与覆盖}\label{ch:matching}
\sourceof{18.315 L7、L24; 18.433 L1--3; Tutte 极大超图证明为编者另证}
\originalsection{增广路}
\begin{definition}
匹配是两两不共端点的边集 $M$. 顶点若关联某条匹配边就称为已匹配, 否则为自由顶点. 匹配覆盖全部顶点时为完美匹配. 匹配数 $\nu(G)$ 是匹配的最大边数. 顶点覆盖是与每条边都至少有一端相交的顶点集, 最小大小记 $\tau(G)$.
\end{definition}
极大匹配不能再加入一条边, 但可能并非最大匹配. 在四顶点路 $1,2,3,4$ 上, 只有中间边的匹配是极大匹配, 而两端边组成大小2的最大匹配. 要改进前者, 必须先删除一条匹配边, 再加入两条边.
\begin{definition}
相对于 $M$, 交替路的边交替属于和不属于 $M$. 增广路是两个端点都是自由顶点、首尾边都不在 $M$ 的交替路. 对增广路翻转边的归属, 得到边数增加1的新匹配.
\end{definition}
\begin{theorem}[Berge 判据]\label{thm:berge}
匹配 $M$ 最大当且仅当不存在相对于它的增广路.
\end{theorem}
\begin{proof}
有增广路就能翻转而增加匹配, 所以非最大. 反向若 $M'$ 更大, 对称差 $M\mathbin\triangle M'$ 的每个顶点度至多2, 因而分支为交替圈、交替路或孤立点. 圈中两种边同样多, 偶长路也同样多. 因总计 $|M'|>|M|$, 至少一路中 $M'$ 边比 $M$ 边多1. 首尾为 $M'$ 边, 端点在 $M$ 中自由, 这就是增广路.
\end{proof}
\begin{center}
\begin{tikzpicture}[x=1.3cm]
\foreach \i/\lab in {0/a,1/u,2/b,3/v,4/c,5/w}\node[vertex](\lab)at(\i,0){$\lab$};
\foreach \u/\v in {a/u,b/v,c/w}\draw[edge,dashed](\u)--(\v);
\foreach \u/\v in {u/b,v/c}\draw[edge,double distance=1.5pt](\u)--(\v);
\end{tikzpicture}
\end{center}
双线为当前匹配, 虚线为非匹配边. 端点 $a,w$ 自由; 翻转整条路后匹配从2条变为3条. 中间顶点仍恰被一条匹配边覆盖.
\originalsection{Hall 定理}
\begin{theorem}[Hall]\label{thm:hall}
二部图 $G=(A\sqcup B,E)$ 存在覆盖全部 $A$ 的匹配, 当且仅当每个 $S\subseteq A$ 都满足 $|N(S)|\ge|S|$.
\end{theorem}
\begin{proof}
必要性: $S$ 中各顶点的匹配伙伴不同且都在 $N(S)$ 中. 充分性对 $|A|$ 归纳. 空侧无须匹配, 单顶点由条件有邻居. 若每个非空真子集 $S\subsetneq A$ 都满足严格不等式, 任取边 $ab$ 并删除其两端; 对剩余侧的非空 $S$, 原邻域至少 $|S|+1$, 删除 $b$ 后仍不少于 $|S|$. 归纳得到剩余匹配, 加回 $ab$ 即可.

若有非空真子集 $S$ 满足 $|N(S)|=|S|$, 在 $S\sqcup N(S)$ 上条件仍成立, 先归纳匹配这部分. 对余图, 任取 $T\subseteq A\setminus S$, 原条件用于 $S\cup T$ 得
\[
|N(S)\cup N(T)|\ge |S|+|T|,
\]
减去 $|N(S)|=|S|$ 得 $|N(T)\setminus N(S)|\ge|T|$. 余图也满足条件, 再匹配后合并两个匹配.
\end{proof}
Hall 条件一次检查所有子集, 并不是说逐个顶点有邻居就够. 三个工人可选的工作分别为 $\{1,2\}$、$\{1,2\}$、$\{1,2\}$, 每人有两选择, 但三人共同邻域只有两项, 不可能全部安排.
\begin{corollary}\label{cor:regularmatching}
有限 $k$ 正则二部多重图在 $k\ge1$ 时有完美匹配.
\end{corollary}
\begin{proof}
两侧度数和相等, 故 $k|A|=k|B|$, 两侧等大. 对 $S\subseteq A$, 从 $S$ 发出的 $k|S|$ 条边都落入 $N(S)$, 而那里至多承接 $k|N(S)|$ 条边. 所以满足 Hall 条件. 定理证明同样适用于多重边, 因为邻域与顶点匹配问题不因重复边改变.
\end{proof}
\originalsection{Kőnig 定理与算法证书}
\begin{theorem}[Kőnig]\label{thm:konig}
有限二部图满足 $\nu(G)=\tau(G)$.
\end{theorem}
\begin{proof}
任意覆盖至少要从每条匹配边选一个端点, 所以 $\tau\ge\nu$. 取最大匹配 $M$, 从 $A$ 侧的所有自由顶点出发, 沿“非匹配边由 $A$ 到 $B$, 匹配边由 $B$ 到 $A$”访问可达顶点. 记可达部分为 $Z_A,Z_B$. 由 Berge, $Z_B$ 没有自由顶点, 否则找到增广路. 每个 $Z_B$ 顶点的匹配伙伴在 $Z_A$ 中; 每个已匹配的 $Z_A$ 顶点的伙伴也在 $Z_B$ 中, 因为到达它的最后一步必须沿该匹配边.

令 $C=(A\setminus Z_A)\cup Z_B$. 若边 $ab$ 没被覆盖, 就有 $a\in Z_A,b\notin Z_B$. 非匹配边会使 $b$ 可达; 匹配边则由伙伴性质使 $b$ 可达, 都矛盾. 因而 $C$ 是覆盖. 每条匹配边的两端同在可达部分或同在不可达部分, 在 $C$ 中恰选一端; 自由顶点均不在 $C$ 中. 所以 $|C|=|M|$.
\end{proof}
这同时给出最大匹配的可检查证书: 提交一个匹配和同样大小的覆盖, 两者分别给出下界与上界. 二部匹配算法从空匹配开始反复搜索增广路, 至多增广 $\min(|A|,|B|)$ 次. 一次搜索扫描顶点与边, 时间 $O(n+m)$, 因而朴素总时间 $O(n(n+m))$. 非二部图中搜索树可能遇到奇圈, 需用 Edmonds 花收缩算法; 本册不把二部搜索直接套到一般图.
\originalsection{一般图的完美匹配}
\begin{theorem}[Tutte]\label{thm:tutte-matching}
有限简单图 $G$ 有完美匹配当且仅当每个 $S\subseteq V$ 都满足 $o(G-S)\le|S|$, 其中 $o(F)$ 为 $F$ 的奇阶分支数.
\end{theorem}
\begin{proof}
若有完美匹配, $G-S$ 的每个奇阶分支至少有一个顶点要与 $S$ 中顶点匹配; 各分支占用不同顶点, 得必要性. 空图也符合定理.

证明充分性. 条件在 $S=\varnothing$ 时保证各分支阶数为偶数, 故总阶数为偶数. 假设没有完美匹配, 在同一顶点集上加边, 取边数极大的无完美匹配超图 $H$. 对每条缺失边 $e$, 图 $H+e$ 都有使用 $e$ 的完美匹配. 令 $U$ 为 $H$ 中与所有其他点相邻的顶点集.

先证 $H-U$ 的每个分支为团. 否则非完全连通分支含诱导路 $x,y,z$, $xz$ 不存在. 因 $y\notin U$, 有顶点 $w$ 不邻接 $y$; 它不同于 $x,y,z$. 取 $H+xz$ 的完美匹配 $M_1$ 和 $H+yw$ 的完美匹配 $M_2$. 两匹配对称差分解成交替偶圈. 如果新增边位于不同圈, 在 $M_1$ 中翻转包含 $xz$ 的圈, 就得到 $H$ 内的完美匹配, 矛盾.

所以两条新增边同在一交替圈 $C$. 删去它们后, 得两条偶长交替路, 每条连接 $\{x,z\}$ 中一点与 $\{y,w\}$ 中一点. 在每条路上取 $M_1$ 边, 只留下来自 $\{x,z\}$ 的端点; 取 $M_2$ 边, 则只留下来自 $\{y,w\}$ 的端点. 若两路分别连接 $x$ 到 $y$、$z$ 到 $w$, 第一条取 $M_2$, 第二条取 $M_1$, 再加原边 $yz$. 若分别连接 $x$ 到 $w$、$z$ 到 $y$, 第一条取 $M_1$, 第二条取 $M_2$, 再加原边 $xy$. 都将圈上各点用 $H$ 的边匹配, 圈外沿用 $M_1$, 矛盾. 结构结论成立.

设 $H-U$ 的奇阶团数为 $q$. 若 $q\le|U|$, 每个奇团取一点与不同的 $U$ 顶点匹配, 各团剩余偶数点在内部配对. $U$ 剩余点数 $|U|-q$ 为偶数, 因总阶数的奇偶性等于 $|U|+q$; $U$ 自身也是团, 所以能内部配对. 又得到完美匹配, 矛盾. 故 $o(H-U)>|U|$. 由于 $G\subseteq H$, 每个 $H-U$ 的奇分支在 $G-U$ 中细分后至少贡献一个奇分支, 所以 $o(G-U)\ge o(H-U)>|U|$, 与原条件矛盾.
\end{proof}
Tutte 条件揭示了真正障碍是“太多奇分支争夺太少外部匹配伙伴”. 例如六顶点星 $K_{1,5}$, 删中心后有5个奇分支, 大于1, 所以偶数阶并不足以保证完美匹配.
\originalsection{习题与解答}
\begin{exercise}\label{ex:match1}
用 Hall 定理证明: 任意 $n\times n$ 非负矩阵若每行、每列和均为1, 则能选出一个置换 $\pi$ 使每个 $a_{i,\pi(i)}>0$.
\end{exercise}
\begin{exercise}\label{ex:match2}
证明任意极大匹配的大小至少是最大匹配的一半.
\end{exercise}
习题\ref{ex:match1}的解答.
\begin{solution}
把行、列作为两侧顶点, 正元素位置连边. 对行集 $S$, 它们总和为 $|S|$, 所有贡献位于 $N(S)$ 中的列, 每列总和为1, 因而 $|S|\le|N(S)|$. 完美匹配对应所求置换.
\end{solution}
习题\ref{ex:match2}的解答.
\begin{solution}
极大匹配 $M$ 的所有端点构成覆盖: 若存在两端都自由的边就还能加入. 因而每条最大匹配边至少占用其中一个端点, 且不同最大匹配边占用不同端点. 得 $\nu\le2|M|$. 四顶点路上的中间边例子达到等号.
\end{solution}

\originalchapter{顶点染色}\label{ch:coloring}
\sourceof{18.315 L4--5、L8; 删双顶点引理及换色过程为编者补证}
\originalsection{染色数与贪心排序}
\begin{definition}
正常顶点染色是映射 $c:V\to\{1,\ldots,k\}$, 使每条边 $uv$ 满足 $c(u)\ne c(v)$. 最小所需颜色数为染色数 $\chi(G)$. 团是内部两两相邻的顶点集, 最大大小记 $\omega(G)$; 最大独立集大小记 $\alpha(G)$.
\end{definition}
染色把顶点分成若干独立集. 团内各点必须异色, 每个颜色类至多含 $\alpha$ 点, 所以 $\chi\ge\omega$, 且对 $n>0$ 有 $\chi\ge n/\alpha$. 这些界不总精确: 奇圈团数为2, 染色数却为3.
\begin{proposition}\label{prop:greedycolor}
有限简单图满足 $\chi(G)\le\Delta+1$. 更一般地, 若每个非空子图有度数至多 $k$ 的顶点, 则 $\chi(G)\le k+1$.
\end{proposition}
\begin{proof}
依任意顺序染点, 一个顶点的已染邻点至多 $\Delta$ 个, 总有一种未被使用的颜色. 对第二结论, 反复在剩余诱导子图删除一个度至多 $k$ 的顶点, 再按删除顺序的逆序染色. 每点染色时至多有 $k$ 个已染邻点.
\end{proof}
第二种图称为 $k$ 退化图. 它控制的是每个子图的稀疏性, 而不是全图平均度. 一个大图的平均度很小, 仍可能含一个巨大团; 只检查全图不能直接给出退化染色界.
\begin{proposition}
若 $m$ 条边的图需要 $q$ 种颜色, 则 $m\ge\binom q2$, 因而 $\chi(G)\le (1+\sqrt{1+8m})/2$.
\end{proposition}
\begin{proof}
在一个最少颜色的染色中, 任意两个颜色类之间必须有边; 否则可把两类合并而少用一色. 不同颜色类对需要不同的边, 所以至少有 $\binom q2$ 条边. 解二次不等式即得上界.
\end{proof}
\originalsection{Brooks 定理}
\begin{lemma}\label{lem:lowdegreecolor}
连通图 $F$ 的最大度不超过 $D$. 若有顶点 $r$ 的度至多 $D-1$, 则 $F$ 可用 $D$ 色正常染色.
\end{lemma}
\begin{proof}
取以 $r$ 为根的生成树, 按深度从大到小染色, 最后染根. 每个非根顶点至少有一个未染的父顶点, 所以已染邻点至多 $D-1$ 个. 根的总度也至多 $D-1$, 贪心过程能完成.
\end{proof}
\begin{lemma}\label{lem:brooksorder}
若 $G$ 为 $D$ 正则、二连通、非完全的简单图, 且 $D\ge3$, 则有非相邻顶点 $a,b$ 及其公共邻点 $v$, 使 $G-\{a,b\}$ 连通.
\end{lemma}
\begin{proof}
若 $G$ 三连通, 选一条非相邻顶点间的最短路, 它的前三点构成诱导路 $a,v,b$. 三连通性直接保证删两端后仍连通.

否则有二点割 $\{v,c\}$. 图 $H=G-v$ 连通, 而 $c$ 是其割点. 由块树引理\ref{lem:blocktree}, $H$ 有两个不同叶块. 每个叶块的内部点集中必须有 $v$ 的邻点: 若没有, 删去该块唯一的连接割点就会把内部点与 $v$ 分离, 违反 $G$ 二连通. 在这两个内部点集中分别取 $a,b\in N_G(v)$. 它们处于不同块的内部, 不相邻. 删去 $a,b$ 后, 每个叶块的剩余部分仍连到原连接点, 故 $H-a-b$ 连通. 由于 $v$ 原有 $D\ge3$ 个邻点, 删除 $a,b$ 后仍有邻点, 所以 $G-a-b$ 也连通.
\end{proof}
\begin{theorem}[Brooks]\label{thm:brooks}
有限简单连通图若既不是完全图, 也不是奇圈, 则 $\chi(G)\le\Delta(G)$.
\end{theorem}
\begin{proof}
最大度不超过2时, 连通图为路径、圈或单点. 排除完全图和奇圈后是非完全路径或偶圈, 都可用2色染色. 以下设 $D=\Delta\ge3$. 如果图非正则, 用低度引理. 若有割点 $v$, 对 $G-v$ 的每个分支 $C_i$ 考虑连通图 $G[C_i\cup\{v\}]$. 它最大度不超过 $D$, 且 $v$ 在其中的度至多 $D-1$, 因为其他分支至少还占有一个邻点. 各部分都可用 $D$ 色染色, 置换各部分颜色使 $v$ 同色后即可拼接.

剩下正则二连通情形, 用上一引理找 $a,v,b$. 先把不相邻的 $a,b$ 同染颜色1. 在 $G-a-b$ 中取以 $v$ 为根的生成树, 从最深点向根贪心染色. 每个非根点有未染父顶点, 可用 $D$ 色. 最后根 $v$ 的两个邻点 $a,b$ 同色, 全部邻点使用的不同颜色至多 $D-1$, 所以根也有颜色可用.
\end{proof}
定理指出 $\Delta+1$ 的例外恰由完全图和奇圈产生. 它没有说任意给定染色都能不加颜色局部改成最优染色; 证明中构造的顶点顺序是实质部分.
\originalsection{Mycielski 构造}
团数与染色数之间的差距可以任意大, 即使禁止三角形也不能仅凭团数控制染色数.
\begin{theorem}
若无三角形图 $G$ 有 $\chi(G)=k\ge2$, 则有无三角形图 $M(G)$ 满足 $\chi(M(G))=k+1$.
\end{theorem}
\begin{proof}
设原顶点为 $v_1,\ldots,v_n$, 加副本 $u_1,\ldots,u_n$ 及一点 $w$. 保留原边; 对每条原边 $v_iv_j$, 加 $u_iv_j,u_jv_i$; 最后令 $w$ 邻接所有 $u_i$. 副本之间无边, $w$ 不邻接原点. 若三角形含 $w$, 另外两点须为两个副本而不相邻, 不可能. 若含一个副本 $u_i$, 另外两原点与 $v_i$ 形成原三角形, 也不可能; 其余三角形都在原图中. 所以新图无三角形.

沿用原 $k$ 色染色, 给 $u_i$ 与 $v_i$ 同色, 给 $w$ 新色, 得到 $k+1$ 色. 假如只需 $k$ 色, 令 $w$ 的颜色为 $k$, 则全部 $u_i$ 都不取 $k$. 在原图中, 将每个颜色 $k$ 的 $v_i$ 改为 $u_i$ 的颜色. 两个被改点原先同色, 本来不相邻; 被改点 $v_i$ 与未改邻点 $v_j$ 的新色不同, 因为新图有边 $u_iv_j$. 得原图的 $k-1$ 色染色, 矛盾.
\end{proof}
从 $C_5$ 开始反复构造, 可得任意高染色数的无三角形图. 构造和第\ref{ch:ramsey}章的随机方法分别展示确定性与概率性两种存在证明.
\originalsection{染色的重新配置}
\begin{theorem}\label{thm:recolor}
图 $G$ 最大度为 $D$, 若 $k\ge D+2$, 则任意两个正常 $k$ 色染色之间, 可通过每次只改变一个顶点的颜色、全程保持正常染色而互相转换.
\end{theorem}
\begin{proof}
对顶点数归纳. 删除一点 $v$, 归纳把余图从初始染色变为目标染色. 模拟余图每次把一点 $u$ 改为颜色 $b$ 的操作. 若 $u$ 不邻接 $v$ 或 $v$ 当前颜色不是 $b$, 直接执行. 否则先给 $v$ 换成一种颜色, 既不同于所有当前邻点颜色, 又不同于 $b$. 被禁颜色至多 $D+1$ 种, 在 $k\ge D+2$ 时总能找到; 随后执行 $u$ 的变化. 余图全部达到目标后, 把 $v$ 改成目标色, 这时所有邻点已是目标色, 不发生冲突.
\end{proof}
颜色数充足保证状态图连通, 不自动给出快速随机采样或最短换色路径. 18.315 L8 还讨论格点三染色的高度函数及直径; 那个专门模型不在这里被泛化为所有图的同一结论.
\originalsection{习题与解答}
\begin{exercise}\label{ex:col1}
证明任何森林可用两色染色, 并给出一个最大度任意大但染色数为2的连通图.
\end{exercise}
\begin{exercise}\label{ex:col2}
证明若 $G$ 的染色数为 $k\ge1$, 则它含有最小度至少 $k-1$ 的子图.
\end{exercise}
习题\ref{ex:col1}的解答.
\begin{solution}
在每棵树中按根的距离奇偶染色; 森林无奇圈, 也可直接用二部判据. 星 $K_{1,D}$ 在 $D\ge1$ 时最大度 $D$ 而染色数2, 说明 Brooks 给出上界而不是精确公式.
\end{solution}
习题\ref{ex:col2}的解答.
\begin{solution}
选顶点数最少的诱导子图 $H$, 使 $\chi(H)=k$. 若某点度至多 $k-2$, 删它后由最小性可用至多 $k-1$ 色染色, 再给它选未被邻点占用的一色, 得 $H$ 的 $k-1$ 色染色, 矛盾. 所以 $\delta(H)\ge k-1$.
\end{solution}

\originalchapter{边染色}\label{ch:edgecolor}
\sourceof{18.315 L6--7; Vizing 扇形的旋转与换色为编者补证}
\originalsection{匹配分解}
\begin{definition}
正常边染色给每条边分配一种颜色, 使有共同端点的两边异色. 最少颜色数为边染色数 $\chi'(G)$. 每个边颜色类就是一个匹配, 所以边染色等价于把全部边分拆成匹配.
\end{definition}
在排班模型中, 顶点可以表示队伍, 边表示两队之间的比赛. 同一队同一时段最多参加一场, 一种边颜色就表示一个时段. 最大度给出必要时段数 $\chi'\ge\Delta$, 但奇圈会需要额外一色: $\chi'(C_{2r+1})=3$, 而 $\Delta=2$.
\begin{theorem}[二部图边染色]\label{thm:bipartiteedge}
有限二部多重图满足 $\chi'(G)=\Delta(G)$.
\end{theorem}
\begin{proof}
若 $\Delta=0$, 无边而结论成立. 在两侧补孤立顶点使它们大小相等. 令 $D=\Delta$, 每个点的缺额度为 $D-d(v)$. 两侧缺额度总和相等, 因为原两侧度数和都等于边数. 反复连接两侧尚有缺额的点, 允许平行边, 直到变成 $D$ 正则二部多重图. 由推论\ref{cor:regularmatching}有完美匹配. 删去它后图为 $D-1$ 正则, 归纳分拆成 $D$ 个完美匹配. 限制到原边, 得 $D$ 色边染色. 下界来自最大度点.
\end{proof}
这证明了 Hall 与边染色之间的联系: 完美匹配移走每个顶点的一单位度, 因而可以逐层剥离正则二部图. 一般图的奇圈使这层结构失效.
\originalsection{双色链与扇形}
以下只讨论简单图. 在一张正常部分边染色中, 只看颜色 $\alpha,\beta$ 的边, 每个顶点至多关联这两色各一条边, 所以分支为路、偶圈或孤立点. 在某个分支交换两色, 不会破坏正常边染色, 这称为 Kempe 换色.
\begin{lemma}[线性扇形旋转]\label{lem:fan}
设只有边 $xy_0$ 未染色, 邻点 $y_0,y_1,\ldots,y_t$ 互异, 且边 $xy_i$ 的颜色 $c_i$ 在 $y_{i-1}$ 处缺失 ($1\le i\le t$). 对任意前缀 $y_0,\ldots,y_s$, 同时把 $xy_{i+1}$ 的颜色移到 $xy_i$ ($i<s$), 让 $xy_s$ 未染, 仍是正常部分边染色.
\end{lemma}
\begin{proof}
中心 $x$ 处保留原有互异颜色的一部分, 没有重复. 在每个 $y_i$ ($i<s$) 处, 移入颜色原本缺失; 最后一个顶点只删去一种颜色. 其余边不变, 所以无冲突. 若颜色 $\gamma$ 在 $x$ 与 $y_s$ 同时缺失, 就可在旋转后给 $xy_s$ 染 $\gamma$.
\end{proof}
旋转必须以原边颜色同时重新赋值, 或从前到后依次正确移位; 不能误把刚改过的颜色继续当作原色传播.
\originalsection{Vizing 定理}
\begin{theorem}[Vizing]\label{thm:vizing}
有限简单图满足 $\Delta\le\chi'\le\Delta+1$.
\end{theorem}
\begin{proof}
固定 $D+1$ 色, 其中 $D=\Delta$. 对边数归纳, 设除 $xy_0$ 外全部边已正常染色. 每点至少缺一色. 固定 $x$ 处缺色 $\alpha$, 从单点扇形 $y_0$ 开始. 在当前末点 $y_t$ 取缺色 $\beta$. 若 $x$ 也缺 $\beta$, 旋转整个扇形后染末边 $\beta$, 完成. 否则 $x$ 恰有一条 $\beta$ 色边 $xz$. 若 $z$ 不在扇形中, 将它作为新末点继续. 中心邻点有限, 所以终会直接完成或重复.

若重复为 $z=y_j$, 则 $1\le j\le t-1$: 首边未染, 且 $y_t$ 缺 $\beta$, 所以重复边既不是首边也不是末边. 有 $c_j=\beta$, 而 $y_{j-1},y_t$ 都缺 $\beta$. 考察 $\alpha,\beta$ 双色子图中含 $x$ 的分支 $P$. $x$ 缺 $\alpha$ 而有 $\beta$, 所以 $P$ 是以 $x$ 为端点的路. 两个缺 $\beta$ 的点若在此路上, 只能成为另一端点, 因而 $y_{j-1},y_t$ 至多一个属于 $P$. 交换 $P$ 上的两色, 此后 $x$ 缺 $\beta$.

如果 $y_{j-1}\notin P$, 它仍缺 $\beta$. 前缀 $y_0,\ldots,y_{j-1}$ 的中心边原来都不是 $\alpha,\beta$: $\alpha$ 在 $x$ 缺失, 唯一 $\beta$ 边是未纳入该前缀的 $xy_j$. 前缀旋转涉及的其他缺色条件不受双色交换影响. 旋转此前缀, 再给其末边染 $\beta$, 完成.

如果 $y_{j-1}\in P$, 则 $y_t\notin P$. 换色后 $y_t$ 仍缺 $\beta$, $y_{j-1}$ 改为缺 $\alpha$, 而 $xy_j$ 由 $\beta$ 改为 $\alpha$. 因而扇形在这个位置的新条件仍成立: 新边色 $\alpha$ 在前点缺失. 其余中心边都不是这两色, 其缺色条件不变. 所以可旋转整个扇形, 再把末边染 $\beta$. 归纳完成. 下界仍由最大度点的关联边必须异色得到.
\end{proof}
这里没有把“沿双色路换色”作为未经验证的口号. 两种情况分别控制换色是否改动旋转前缀, 这正是原手写稿略写、必须补齐的地方. 简单图条件也很重要: 平行边会让不同中心边共用另一端点, 上述不同邻点扇形构造不能直接使用.
\originalsection{完全图的时段安排}
\begin{proposition}
$\chi'(K_{2r})=2r-1$, $\chi'(K_{2r+1})=2r+1$ ($r\ge1$).
\end{proposition}
\begin{proof}
对 $K_{2r}$, 把顶点标为 $\infty$ 及模 $2r-1$ 的剩余类. 对每个 $a$, 取匹配
\[
\{\infty a\}\ \cup\ \{(a+i)(a-i):1\le i\le r-1\}.
\]
模数为奇数, 非零剩余类可按 $\pm i$ 成对, 所以每个匹配覆盖全部顶点. 任意两有限标号 $u,v$ 的边归于唯一中点 $a=(u+v)/2$, 除2在奇模数下可逆; $\infty$ 的边也各出现一次. 得到 $2r-1$ 色, 与最大度下界相等.

$K_{2r+1}$ 每个匹配最多含 $r$ 条边, 而总边数为 $r(2r+1)$, 所以至少需 $2r+1$ 色. 把它嵌入 $K_{2r+2}$, 用刚才的 $2r+1$ 色安排再删除额外顶点即可.
\end{proof}
\originalsection{习题与解答}
\begin{exercise}\label{ex:edge1}
证明图的边染色数等于其线图的顶点染色数. 线图的顶点为原图各边, 共享端点的边在新图中相邻.
\end{exercise}
\begin{exercise}\label{ex:edge2}
一所学校有若干老师与班级, 每条教学任务连接一位老师和一个班级. 每位老师及每个班级各参与至多 $D$ 条任务, 同一老师与班级可有多条任务. 证明能用 $D$ 个时段安排全部任务, 每时段老师和班级都没有冲突.
\end{exercise}
习题\ref{ex:edge1}的解答.
\begin{solution}
线图每个顶点颜色对应原边颜色. 线图相邻顶点异色恰好要求原图共端点两边异色, 两种染色逐一对应. 但在线图套一般 $\Delta+1$ 贪心界通常不如 Vizing, 因为原边 $uv$ 在线图度数最多 $d(u)+d(v)-2$.
\end{solution}
\Needspace{4\baselineskip}
习题\ref{ex:edge2}的解答.
\begin{solution}
任务构成二部多重图, 两侧为老师和班级. 二部图边染色定理给至多 $D$ 色, 每个颜色类为不共端点的匹配, 对应一个可同时执行的时段. 无需把重复任务合并, 因为多重图版本已保留它们.
\end{solution}

\originalchapter{平面图与对偶}\label{ch:planar}
\sourceof{6.042J 第12章; 18.315 L6--7、L10; 桥与面边界计数补全}
\originalsection{嵌入与面}
\begin{definition}
平面图是一张已经在平面中画出的图, 边的内部彼此不交且不经过其他顶点. 可平面图是存在这种画法的抽象图. 平面去掉所画顶点与边后的连通区域称为面, 包括唯一无界外面. 面的度为沿其全部边界行走得到的边出现总次数.
\end{definition}
同一抽象图可有不同嵌入, 但连通图的面数由顶点与边数决定. 桥两侧属于同一个面, 在面边界中计两次. 因而即使一个面不是由简单圈围成, 仍有面度总和 $\sum_F d(F)=2m$.

以下采用平面拓扑中的 Jordan 圈分离定理作为外部先修工具: 一条简单闭曲线把球面分成两个连通区域. 本书不证明该拓扑定理, 而是明确用它支持“圈边两侧为不同面”以及五色证明中的不能交叉. 平面加上无穷远点即为球面, 外面与其他面于是平等.
\begin{theorem}[Euler 公式]\label{thm:eulerformula}
平面嵌入图有 $n$ 个顶点、$m$ 条边、$f$ 个面和 $c$ 个分支时, 满足 $n-m+f=1+c$.
\end{theorem}
\begin{proof}
连通图若无圈, 是树, 有 $m=n-1$, 且补集只有一个面, 公式成立. 若有圈, 删除圈上的一边. 两侧面不同, 删除边将它们合并, 所以边数和面数各减1; 图仍连通. 反复删除直到树, $n-m+f$ 始终不变, 因而为2.

非连通图可在同一个面的不同分支之间加边, 连接 $c-1$ 次后连通. 新边是桥, 不分割面, 面数不变. 对新图用连通公式: $n-(m+c-1)+f=2$, 整理得到结论.
\end{proof}
\begin{corollary}\label{cor:planaredges}
简单可平面图若 $n\ge3$, 则 $m\le3n-6$. 若还无三角形, 则 $m\le2n-4$. 因而 $K_5$ 和 $K_{3,3}$ 都不可平面.
\end{corollary}
\begin{proof}
先设连通. 当 $n\ge3$ 且简单时, 每个面边界长度至少3; 长度1需要自环, 长度2需要平行边或整图只是一条边. 所以 $3f\le2m$, 联合 $f=2-n+m$ 得第一界. 无三角形时每面边界长度至少4, 得第二界. 树的边界可能重复边, 但在 $n\ge3$ 时长度 $2(n-1)\ge4$, 与计算一致.

非连通图在面内加桥使连通, 不产生圈, 因而也不产生三角形, 再应用两界. $K_5$ 有10边而 $3n-6=9$; $K_{3,3}$ 二部而无三角形, 有9边而 $2n-4=8$.
\end{proof}
少于界限也不保证可平面. 例如把 $K_{3,3}$ 每条边细分一次, 边数与顶点数增加后仍不可平面, 但可以满足 $m\le3n-6$. 细分不改变是否能无交叉画图.
\originalsection{五色定理}
\begin{theorem}[五色定理]\label{thm:fivecolor}
每个有限简单可平面图都可正常地用五色染色.
\end{theorem}
\begin{proof}
对顶点数归纳. 平面边数界与握手引理保证至少一点 $v$ 度至多5. 删除它后归纳染色. 若邻点不足5个或所用颜色不足5种, 给 $v$ 选未用色即可. 否则按嵌入中围绕 $v$ 的循环顺序, 五邻点为 $v_1,\ldots,v_5$, 其颜色依次为1至5.

考虑仅有1、3两色的诱导子图. 若 $v_1,v_3$ 不在同一分支, 交换含 $v_1$ 分支的1、3两色, 保持正常染色并在 $v$ 邻域释放颜色1. 若它们相连, 取一条简单双色路, 加上 $vv_1,vv_3$ 形成简单闭曲线. 邻点 $v_2,v_4$ 位于其不同侧. 2、4双色路不能经过1、3双色路的顶点, 也不能穿越其边, 因此 $v_2,v_4$ 在2、4诱导子图中不连通. 交换含 $v_2$ 分支的2、4两色, 释放颜色2, 再给 $v$ 染色.
\end{proof}
四色定理进一步把五色改为四色, 但其证明不属于本册. 五色定理不是四色定理的证明, 两条双色链能否同时交换必须按平面分离关系检查.
\begin{center}
\begin{tikzpicture}[scale=1.05]
\node[vertex](v)at(0,0){$v$};
\foreach \i/\ang in {1/90,2/18,3/-54,4/-126,5/162}{\node[vertex](v\i)at(\ang:1.8){$v_{\i}$};\draw[edge](v)--(v\i);}
\draw[very thick,color=structurecolor](v1)..controls(3.4,2.6)and(3.4,-2.3)..node[pos=.48,right=3mm]{$1,3$路}(v3);
\node at(0,-2.5){循环顺序决定 $v_2,v_4$ 分处两侧};
\end{tikzpicture}
\end{center}
\originalsection{平面对偶}
\begin{definition}
连通平面图 $G$ 的对偶图 $G^*$ 在每个面中放一个顶点, 对每条原边加一条穿过它一次的对偶边, 连接其两侧的面. 对偶通常是允许自环和平行边的多重图. 原边与对偶边一一对应.
\end{definition}
\begin{proposition}\label{prop:dual-tree}
若 $T$ 为 $G$ 的生成树, 则原图不在 $T$ 中的边对应的对偶边组成 $G^*$ 的生成树. 原图的桥对应对偶自环.
\end{proposition}
\begin{proof}
桥两侧同一面, 因而对应自环. 对生成树, 先只画树, 补集连通. 逐条加入原图的非树边; 每次新边都与现有图中的路构成圈, 把一个面分成两个. 在追踪面分裂的对偶记录中, 新边连接这两个子面, 相当于把一个已有顶点分裂成两个并以新边连接. 从单点开始这样构造, 所得记录始终是一棵树. 最终它恰为全部非树边的对偶边, 连接全部 $f$ 个面, 边数为 $m-(n-1)=f-1$.
\end{proof}
\originalsection{曲面上的上界}
平面是球面去掉一点. 若允许把图画在有 $g$ 个把手的闭可定向曲面上, 圈分割关系就不同. 只在这里补充原课中的上界, 不把完整曲面分类或 Heawood 最优性证明纳入核心.
\begin{proposition}
若简单图可嵌入亏格 $g\ge1$ 的闭可定向曲面, 则
\[
\chi(G)\le H(g):=\left\lfloor\frac{7+\sqrt{1+48g}}2\right\rfloor.
\]
\end{proposition}
\begin{proof}
所用曲面 Euler 工具为 $n-m+f=2-2h$, 适用于每个面为圆盘的连通嵌入. 一般嵌入可取图的窄带邻域, 给每条边界圈封圆盘, 得亏格 $h\le g$ 的这种嵌入. 面边界长度至少3时得到 $m\le3n-6+6g$; 对 $n\ge3$ 的简单连通图仍适用, 与平面计数同理.

若染色数 $k\le7$, 因 $H(g)\ge7$ 已有结论. 若 $k\ge8$, 取第\ref{ch:coloring}章习题中的 $k$ 临界子图, 它有 $n\ge k$ 点且最小度至少 $k-1$, 并且连通: 否则某个较小分支已具有染色数 $k$, 与临界选择矛盾. 嵌入限制在子图上仍可用上述边界, 故 $(k-1)n\le6n-12+12g$, 即 $(k-7)n\le12g-12$. 因 $k-7>0$, 有 $k(k-7)\le12g-12$, 解二次式得 $k\le(7+\sqrt{1+48g})/2$, 再取整.
\end{proof}
曲面 Euler 公式和“封边界不增加亏格”在这里作为明确的拓扑工具引用. 球面 $g=0$ 的四色结论不由这个推导给出. Kuratowski 的两个细分禁图判据以及 Tutte 的四连通平面图 Hamilton 定理也只列为进一步阅读, 不作为正文后续推理的前提.
\originalsection{习题与解答}
\begin{exercise}\label{ex:plan1}
若连通简单平面图的每个面都为三角形, 求边数与面数.
\end{exercise}
\begin{exercise}\label{ex:plan2}
证明每个简单二部可平面图都有度数至多3的顶点, 并由此给出四色染色算法.
\end{exercise}
习题\ref{ex:plan1}的解答.
\begin{solution}
面度和为 $3f=2m$, Euler 给 $n-m+f=2$, 得 $m=3n-6$, $f=2n-4$. 条件包括外面也为三角形, 不可漏算它.
\end{solution}
习题\ref{ex:plan2}的解答.
\begin{solution}
$n\ge3$ 时无三角形边界给平均度 $2m/n\le4-8/n<4$, 所以有度至多3点; 小图直接成立. 每个诱导子图仍二部且可平面, 可不断删度至多3点, 再逆序贪心用四色. 实际由二部性只需两色; 此处算法用于说明边数界如何给退化上界.
\end{solution}

\originalchapter{染色与 Tutte 多项式}\label{ch:polynomials}
\sourceof{18.315 L9--12; 计数、符号消去和活动性证明补全}
\originalsection{删除与收缩}
\begin{definition}
对非负整数 $q$, 令 $P_G(q)$ 为图 $G$ 的正常 $q$ 色顶点染色数. 缩边 $G/e$ 把非自环边 $e=uv$ 的两端识别为一点, 其余边随端点一起映射; 允许出现平行边和自环. 平行边不增加顶点染色约束, 自环使正常染色数为0.
\end{definition}
\begin{theorem}\label{thm:chromaticdel}
对非自环边 $e$, $P_G(q)=P_{G-e}(q)-P_{G/e}(q)$. 因而 $P_G$ 是多项式; 无自环 $n$ 顶点图的首项为 $q^n$.
\end{theorem}
\begin{proof}
删边图的染色分为端点异色和同色两类. 前者是原图染色; 后者与缩边图的染色一一对应. 无边图有 $q^n$ 种染色. 对边数归纳, 删除项顶点数不变, 收缩项少一个顶点, 得多项式与首项. 对平行边可先删除重复约束, 不影响结论.
\end{proof}
树的染色多项式为 $q(q-1)^{n-1}$: 根有 $q$ 种选择, 按父先子后, 每点只需避开父色. 完全图有 $q(q-1)\cdots(q-n+1)$. 对 $n\ge3$ 的圈, 删除一边得到树, 收缩得到更短的圈; $C_3$ 可直接计算, 归纳得
\[
P_{C_n}(q)=(q-1)^n+(-1)^n(q-1).
\]
这是计数恒等式, 因在所有非负整数成立, 也作为多项式恒等式在任意实数或复数成立.
\originalsection{容斥与断圈}
令 $c(A)$ 为生成子图 $(V,A)$ 的分支数, 包括孤立点. 对每条边的坏事件“端点同色”做容斥, 同时要求 $A$ 中边端点同色时, 每个分支只能选一个公共颜色, 得
\begin{equation}\label{eq:chromaticsubset}
P_G(q)=\sum_{A\subseteq E}(-1)^{|A|}q^{c(A)}.
\end{equation}
\begin{theorem}[Whitney 断圈定理]\label{thm:nbc}
给简单图的边固定全序. 一个断圈是从某个圈删去其最大边所得的边集. 令 $b_i$ 为不包含任何断圈的 $i$ 边子集数, 则 $P_G(q)=\sum_i(-1)^i b_i q^{n-i}$. 特别地, 系数按符号交替, $b_0=1$, $b_1=m$.
\end{theorem}
\begin{proof}
在容斥和中, 若 $A$ 含断圈, 存在一条边 $e$ 的两端可由 $A$ 中小于 $e$ 的边连接. 称其为可切换边, 选最大的这种 $e$, 将 $A$ 改为 $A\mathbin\triangle\{e\}$. 加删 $e$ 都不改变分支数, 而使符号相反.

须检查选择没有改变. 小于或等于 $e$ 的边的可切换性不依赖 $e$ 本身. 对大于 $e$ 的候选边 $f$, 若一条连接路使用 $e$, 可用原先小于 $e$ 的连接路替换它, 所以加入或删除 $e$ 不改变由小于 $f$ 边产生的可达关系. 因而最大可切换边仍为 $e$, 操作是无固定点的对合, 这些项两两消去.

剩余边集不含断圈, 特别不能含圈, 因而为森林, $c(A)=n-|A|$. 按边数归类得公式. 简单图中没有一边断圈, 所以每条单边集都计入 $b_1$.
\end{proof}
原手写稿把断圈定义为删去最小边; 将全序反转就与这里的版本一致. 两种约定不能在同一证明中混用.
\originalsection{无圈定向}
\begin{theorem}[Stanley 的特殊取值]\label{thm:stanley}
无自环图 $G$ 的无有向圈定向数 $a(G)$ 满足 $a(G)=(-1)^nP_G(-1)$.
\end{theorem}
\begin{proof}
对任一非自环边 $e=uv$, 考察 $G-e$ 的无圈定向. 若已有从 $u$ 到 $v$ 的有向路, 新边只能定向为 $u\to v$; 反方向会形成圈. 若已有反向路则同理. 无圈性排除两种路同时存在. 若两方向均不可达, 新边有两种选择.

这种两端不可达的定向与 $G/e$ 的无圈定向对应: 识别 $u,v$ 后不会新增有向圈, 否则识别前有一条两端连接路. 若有共同邻点, 两条边的方向必须一致指入或指出, 否则产生 $u$--$v$ 有向路, 所以重复边的方向相容. 反过来展开识别点也得到唯一的这种定向. 因而 $a(G)=a(G-e)+a(G/e)$. 有向多重图中一对端点相同、方向相反的弧组成长度2的有向圈. 因而平行边在无圈定向中必须同向, 可合并; 自环使 $a=0$. 无边图有唯一空定向. 将染色删除收缩式乘以 $(-1)^n$, 因收缩少一点, 得相同的加法递推与基值, 归纳即得.
\end{proof}
\begin{example}
求 $C_4$ 的染色多项式与无圈定向数.
\end{example}
\begin{solution}
$P_{C_4}(q)=(q-1)^4+(q-1)=q^4-4q^3+6q^2-3q$. 代入 $q=-1$ 得14. 也可直接看16种边方向中, 恰有顺时针与逆时针两种形成有向圈, 所以剩14种.
\end{solution}
\originalsection{Tutte 多项式}
顶点染色只是许多边子集计数之一. 为同时记录“还欠多少连接”和“已经多出多少圈”, 定义秩 $r(A)=n-c(A)$, 零度 $|A|-r(A)$. 二者分别是该生成子图的生成森林边数和多余边数.
\begin{definition}
允许自环和平行边的有限图的 Tutte 多项式为
\[
T_G(x,y)=\sum_{A\subseteq E}(x-1)^{r(E)-r(A)}(y-1)^{|A|-r(A)}.
\]
两个指数都是非负整数, 无边图的值为1.
\end{definition}
\begin{theorem}\label{thm:tuttedel}
普通非桥非自环边满足 $T_G=T_{G-e}+T_{G/e}$. 桥满足 $T_G=xT_{G/e}$, 自环满足 $T_G=yT_{G-e}$.
\end{theorem}
\begin{proof}
把子集分为含 $e$ 和不含 $e$ 两类. 对普通边, 删边不改变 $r(E)$; 不含边的部分直接为删除项. 含边写为 $B\cup\{e\}$, 缩边后秩变为 $r_G(B\cup\{e\})-1$, 全图秩也减1, 零度不变, 得收缩项.

若 $e$ 为桥, 令 $H=G/e$. 对 $B\subseteq E\setminus\{e\}$, 有 $r_G(B)=r_H(B)$、$r_G(B\cup\{e\})=r_H(B)+1$, 全图秩则满足 $r_G(E)=r_H(E(H))+1$. 所以不含桥的项多出 $(x-1)$, 含桥的项对应因子1, 两者相加为 $x$. 若 $e$ 为自环, 加它不改变秩, 却使零度加1, 因子为 $1+(y-1)=y$.
\end{proof}
\begin{proposition}\label{prop:tuttevalues}
设 $G$ 连通. $T_G(1,1)$ 是生成树数, $T_G(2,1)$ 是森林边子集数, $T_G(1,2)$ 是连通生成子图数, $T_G(2,2)=2^m$. 还满足
\[
P_G(q)=(-1)^{n-c(G)}q^{c(G)}T_G(1-q,0),\qquad a(G)=T_G(2,0).
\]
\end{proposition}
\begin{proof}
在 $(1,1)$ 处, 只有两个指数都为0的子集保留, 即连通且无圈的生成树; $0^0$ 在多项式代入时按常数1理解. 在 $(2,1)$ 处只保留零度0的森林; 在 $(1,2)$ 处只保留秩达到全秩的子集; 在 $(2,2)$ 处每子集贡献1.

代入 $(1-q,0)$ 后, 子集项为 $(-q)^{r(E)-r(A)}(-1)^{|A|-r(A)}$. 乘前置因子使总符号为 $(-1)^{|A|}$, $q$ 指数为 $n-r(A)=c(A)$, 与容斥式一致. 再用 Stanley 公式取 $q=-1$, 得 $a=T_G(2,0)$.
\end{proof}
\originalsection{生成树的活动性}
\begin{definition}
固定连通图边的全序及生成树 $T$. 非树边 $e$ 加到树中形成的唯一圈为基本圈; 若 $e$ 为其中最大边, 称它外部活动. 树边 $f$ 删去后产生的两部分之间所有原边构成基本割; 若 $f$ 为其中最大边, 称它内部活动. 记活动边数为 $e(T),i(T)$.
\end{definition}
\begin{theorem}\label{thm:activities}
$T_G(x,y)=\sum_{T\text{ 生成树}}x^{i(T)}y^{e(T)}$.
\end{theorem}
\begin{proof}
对边数归纳, 每次选全序中最小边 $h$. 若它是桥, 每棵生成树都含它, 基本割只有自身, 因而内部活动; 缩去它后其他活动判定不变, 得因子 $x$. 若是自环, 每棵树都不含它, 基本圈只有自身, 外部活动, 得因子 $y$.

若 $h$ 普通, 它在任何生成树中都不活动: 含它时基本割有另一边, 不含时基本圈有另一边, 而 $h$ 是最小的. 不含 $h$ 的生成树对应 $G-h$ 的生成树; 含 $h$ 的对应 $G/h$ 的生成树. 对其他边, 基本圈或基本割在删除收缩下仅可能丢去最小边 $h$, 或将经过 $h$ 的连通结构识别, 最大边身份不改变. 具体地, 树中其他边的基本割对应切开同样的两部分, 非树边的基本圈对应收缩同一路段; 删除项只是从基本割候选中去掉 $h$. 因而其他边活动性全保留. 两类求和与 Tutte 普通边递推相同. 无边单顶点为基值1, 归纳完成.
\end{proof}
这也解释系数为什么非负, 尽管子集表达式中有 $(x-1),(y-1)$. 活动性依赖所选全序, 而整个和与全序无关. 对圈 $C_n$, 递推给 $T_{C_n}=x^{n-1}+x^{n-2}+\cdots+x+y$, 在 $(1,1)$ 处得到 $n$ 棵生成树.

对连通平面图还有 $T_{G^*}(x,y)=T_G(y,x)$. 证明可直接用删缩递推: 原边的删除对应对偶边收缩, 原边收缩对应对偶边删除, 桥与自环交换. 桥与自环用交换后的乘法递推, 不对删桥后的非连通图调用这里的对偶定义. 单顶点无边图及其对偶的值均为1, 归纳即得. 因而树数在对偶下相等, 与命题\ref{prop:dual-tree}的双射一致.
\originalsection{习题与解答}
\begin{exercise}\label{ex:poly1}
求含 $n$ 点、$c$ 个分支的森林的染色多项式和 Tutte 多项式.
\end{exercise}
\begin{exercise}\label{ex:poly2}
证明 $n$ 个标号顶点上的连通简单图数 $c_n$ 满足 $c_1=1$ 及
\[
c_n=2^{\binom n2}-\sum_{j=1}^{n-1}\binom{n-1}{j-1}c_j2^{\binom{n-j}2}.
\]
\end{exercise}
\Needspace{4\baselineskip}
习题\ref{ex:poly1}的解答.
\begin{solution}
各树独立染色, 每个根有 $q$ 种选择, 其余点有 $q-1$ 种, 得 $q^c(q-1)^{n-c}$. 全部 $n-c$ 条边为桥, 所以 $T=x^{n-c}$. 孤立点不增加 Tutte 因子, 但增加染色多项式的 $q$ 因子.
\end{solution}
习题\ref{ex:poly2}的解答.
\begin{solution}
全部图数为 $2^{\binom n2}$. 按包含标号1的连通分支的大小 $j$ 分类: 另选 $j-1$ 点有 $\binom{n-1}{j-1}$ 种, 内部为一个连通图有 $c_j$ 种, 其余点任意成图有 $2^{\binom{n-j}2}$ 种, 两部分之间没有边. 对 $j=1,\ldots,n$ 求和, 把 $j=n$ 项移到左边. 这是原课完全图 Tutte 取值的一个可独立复算的计数实例.
\end{solution}

\originalchapter{搜索与最短路}\label{ch:shortest}
\sourceof{6.042J 第9章有向图与调度部分; 本章算法推导为编者补充}
\originalsection{表示与搜索}
图算法的运行时间依赖输入如何存储. 邻接矩阵是 $n\times n$ 矩阵, 第 $uv$ 项记录是否有边或其权. 它可在常数时间查询一条边, 但占 $O(n^2)$ 空间. 邻接表为每点列出邻边, 总大小 $O(n+m)$; 搜索稀疏图通常采用邻接表. 多重边需分配独立边编号, 不能只用端点识别.

广度优先搜索 BFS 从起点 $s$ 出发, 使用先进先出队列. 初始化 $d(s)=0$, 其他距离为未知. 取出队首 $u$, 扫描其每个邻点 $v$; 若 $v$ 未访问, 设置 $d(v)=d(u)+1$, 父顶点为 $u$, 并入队. 每点最多入队一次, 每条无向边扫描两次, 有向弧扫描一次, 总时间 $O(n+m)$.
\begin{theorem}\label{thm:bfs}
BFS 首次访问顶点时记录的距离等于无权图最短路长度. 父边在可达部分构成最短路树.
\end{theorem}
\begin{proof}
队列按记录距离非递减取出: 距离 $k$ 点新加入的邻点全为 $k+1$, 它们排在尚存距离 $k$ 点之后. 父边给出一条长度 $d(v)$ 的实际路, 所以真实距离不大于记录值. 反向按真实距离归纳, 若最短路中 $v$ 的前点 $u$ 距离为 $k-1$, 它先被取出时将 $v$ 至迟以距离 $k$ 访问. 因而记录值不大于真实距离. 父顶点距离严格减1, 所以无圈, 沿父边都回到根.
\end{proof}
深度优先搜索 DFS 则沿尚未访问的邻点递归, 无新邻点时回退. 顶点分为白色未访问、灰色在递归栈中、黑色已完成. 每点进入和退出各一次, 也需 $O(n+m)$ 时间. 在有向图中, 指向灰色顶点的弧与递归栈上的路组成有向圈; 有向圈在 DFS 中必产生这样的弧, 因为从圈上最先访问点继续追踪时, 在它退出之前圈上可达点都必须被处理.
\originalsection{有向无圈图与调度}
\begin{definition}
有向无圈图简称 DAG. 拓扑排序是顶点的线性顺序, 使每条弧都由较前顶点指向较后顶点.
\end{definition}
\begin{theorem}\label{thm:toposort}
有限有向图有拓扑排序当且仅当它是 DAG. 反复删除入度0顶点即可构造排序.
\end{theorem}
\begin{proof}
拓扑顺序沿弧严格前进, 不可能形成有向圈. 若 DAG 没有入度0顶点, 从任一点不断沿入弧倒走, 有限顶点保证重复, 产生有向圈, 矛盾. 删除一个入度0点后仍无圈, 归纳排序余图, 把删点放在最前即可. 用队列维护入度, 每删除点时减少其后继的入度, 每弧处理一次, 时间 $O(n+m)$.
\end{proof}
若每项工作 $v$ 耗时 $p(v)\ge0$, 弧 $u\to v$ 表示 $v$ 必须等 $u$ 完成, 无限并行资源下最早完工时间满足 $F(v)=p(v)+\max_{u\to v}F(u)$, 没有前驱时最大值记0. 按拓扑顺序计算, 每个前驱值都已确定. 任意调度的完成时间至少达到这个值, 因最长先后依赖链必须串行; 按这些最早开始时间安排又满足全部约束, 所以这是可实现的最优工期. 资源人数限制会改变问题, 不能直接沿用此结论.
\originalsection{松弛与 Dijkstra}
对弧 $uv$ 赋实权 $w(uv)$, 路权为边权和. 若起点可达负权圈, 经过它重复行走可使某些终点路权趋于 $-\infty$, 此时不存在有限最短行走. Dijkstra 假设全部边权非负.

距离松弛是把估计值 $D(v)$ 改为 $\min\{D(v),D(u)+w(uv)\}$. 从 $D(s)=0$, 其余为 $+\infty$ 出发, 每个有限估计都对应实际行走, 因而它始终是最短距离的上界. Dijkstra 反复从未确定点中选最小估计点 $u$, 将它确定, 再松弛出弧.
\begin{theorem}\label{thm:dijkstra}
非负权图中, Dijkstra 确定 $u$ 时有 $D(u)=\dist_w(s,u)$.
\end{theorem}
\begin{proof}
假设已有确定点的值均正确. 对新选 $u$, 取一条最短路; 非负权保证可删圈而得到简单最短路. 若 $u$ 不可达, 所有未确定点估计均无限, 算法可以结束. 否则令 $y$ 为该路上第一个未确定点, 前点 $x$ 已确定, 松弛后 $D(y)\le\dist_w(s,y)$. 结合上界不变量, 等号成立. 非负性保证路径后段权非负, 所以 $\dist_w(s,y)\le\dist_w(s,u)$. 选择最小估计给 $D(u)\le D(y)\le\dist_w(s,u)$, 与始终为上界合起来得等号. 起点为归纳基点.
\end{proof}
用矩阵或顺序扫描选最小值, 时间 $O(n^2+m)$. 用邻接表与带顶点位置索引的二叉堆, 每点在堆中至多一个项, 成功松弛做减小键值, 时间 $O((n+m)\log n)$. 若改为加入新项并跳过过期项, 堆最多有 $O(n+m)$ 项, 对任意多重图应写 $O((n+m)\log(n+m))$; 在简单图 $m\le n(n-1)$ 时仍可写前一界. 算法需要的非负条件是正确性条件, 不只是程序优化建议.
\begin{example}
有向弧为 $sa:4$, $sb:1$, $ba:2$, $ac:1$, $bc:5$, $ct:3$, $at:7$. 从 $s$ 求最短距离.
\end{example}
\begin{solution}
先确定 $s$, 得 $D(a)=4,D(b)=1$. 确定 $b$ 后, $a$ 改为3, $c$ 为6. 确定 $a$ 后, $c$ 改为4, $t$ 为10. 确定 $c$ 后, $t$ 改为7. 最终顺序 $s,b,a,c,t$, 距离 $0,1,3,4,7$. 最短路 $s,b,a,c,t$ 的权为 $1+2+1+3=7$, 父边同时提供可验证路径.
\end{solution}
\originalsection{负权与 Bellman--Ford}
\begin{theorem}\label{thm:bellmanford}
定义 $D_k(v)$ 为从 $s$ 到 $v$ 使用至多 $k$ 条弧的最小行走权, 不存在时为无限. 则
\[
D_{k+1}(v)=\min\bigl\{D_k(v),\min_{uv\in E}(D_k(u)+w(uv))\bigr\}.
\]
若没有从 $s$ 可达的负圈, $D_{n-1}$ 为最短距离. 若 $D_n(v)<D_{n-1}(v)$ 对某点成立, 则有可达负圈.
\end{theorem}
\begin{proof}
至多 $k+1$ 弧的行走或者原已至多 $k$ 弧, 或最后一步为 $uv$, 之前至多 $k$ 弧, 得递推. 无可达负圈时可删除行走中的非负圈, 最短值由至多 $n-1$ 弧的简单路实现. 若第 $n$ 轮仍严格改进, 一条改进的行走有 $n$ 条弧而至少重复一点. 若所含所有圈非负, 删圈得到至多 $n-1$ 弧且权不增加的行走, 与严格改进矛盾; 因而含可达负圈.
\end{proof}
上式用两份数组保存相邻轮, 是最容易核查的不变量版本, 时间 $O(nm+n^2)$, 空间 $O(n)$ 可再配父指针. 通常把 $n-1$ 轮所有边松弛称为 Bellman--Ford; 原地更新可以在一轮传播多条边, 但仍保证最终结果, 因为它不会增加估计且至少包含上述一轮的改进. 第 $n$ 轮能够改进的点及其后继, 最短行走值为 $-\infty$; 与负圈无关的点仍可能有有限最短距离.
\originalsection{全源最短路}
Floyd--Warshall 令 $D^{(k)}_{ij}$ 为内部顶点只允许来自 $\{1,\ldots,k\}$ 的最短路权. 非对角初值为直接弧最小权, 无弧时为无限; 对角初值为0与该点所有自环权的最小值, 没有自环时为0. 这样不会漏掉负权自环. 在无负圈条件下,
\[
D^{(k)}_{ij}=\min\{D^{(k-1)}_{ij},D^{(k-1)}_{ik}+D^{(k-1)}_{kj}\}.
\]
一条简单最短路或不经过 $k$, 或在 $k$ 分成两段, 所以递推正确. 三重循环时间 $O(n^3)$, 矩阵空间 $O(n^2)$. 若最终某个对角值为负, 图有负圈; 必须报告这个条件, 不能把负对角仍称为正常有限最短路矩阵.
\originalsection{习题与解答}
\begin{exercise}\label{ex:short1}
构造一张有负权弧而无负圈的图, 使 Dijkstra 的确定步骤出错.
\end{exercise}
\begin{exercise}\label{ex:short2}
无向图每边权为1. 解释为何 BFS 比二叉堆 Dijkstra 更直接, 并说明权只取0、1时可怎样用双端队列.
\end{exercise}
习题\ref{ex:short1}的解答.
\begin{solution}
取弧 $sa:2,sb:5,ba:-4$, 图无有向圈. Dijkstra 先确定 $a$ 为2, 但真正距离沿 $s,b,a$ 为1. 负权使最短路上较后点距离反而更小, 破坏证明中的非递减关系.
\end{solution}
习题\ref{ex:short2}的解答.
\begin{solution}
单位权下新点距离恰比出队点多1, 普通队列自动保持层序, 时间 $O(n+m)$. 0、1权下每次成功松弛权0边将邻点放队首, 权1边放队尾, 保存过期距离用于跳过旧项, 就维持队列按距离层次排列. 从最小层向外处理的正确性与非负权证明一致, 这是0--1 BFS; 用标准的确定标志实现时每点仅在首次以最小距离出队时扫描邻边, 时间仍为 $O(n+m)$.
\end{solution}

\originalchapter{最小生成树与最小割}\label{ch:mst}
\sourceof{18.433 L9 的随机缩边; 生成树交换与贪心算法为编者补充}
\originalsection{割性质与圈性质}
设连通无向图每边有实权 $w(e)$. 最小生成树是总权最小的生成树. 图有限, 树的数量有限且至少一棵, 所以最小树存在; 边权可为负, 不需要最短路算法的非负假设.
\begin{theorem}[安全边]\label{thm:safeedge}
设边集 $A$ 已包含于某棵最小生成树. 若一个割不被 $A$ 中任何边穿过, 则该割中的一条最轻边 $e$ 可加入 $A$, 新集合仍包含于某棵最小生成树.
\end{theorem}
\begin{proof}
取包含 $A$ 的最小树 $T$. 若 $e\in T$ 已完成. 否则加 $e$ 后产生唯一圈, 圈沿 $e$ 跨割后必须有另一条边 $f$ 跨回. 因割不穿过 $A$, $f\notin A$. 又 $w(e)\le w(f)$, 用 $e$ 替换 $f$ 得一棵包含 $A\cup\{e\}$ 的树, 权不增加. 由原树最优性, 新树也是最小树.
\end{proof}
\begin{proposition}[圈排除]
若边 $e$ 在某个圈上严格比其他边重, 它不在任何最小生成树中. 若各边权互异, 最小生成树唯一.
\end{proposition}
\begin{proof}
树若含 $e$, 删它分两块, 原圈上必有另一边 $f$ 跨两块. 替换后权严格减少, 矛盾. 唯一性方面, 若两棵最小树 $T,T'$ 不同, 取其对称差中权最小的边 $e$, 可设它在 $T$ 中. 删除 $e$ 得割, $T'$ 中跨割的路含某条 $f\notin T$, 否则原 $T-e$ 已跨割. 最小性与权互异使 $w(e)<w(f)$, 在 $T'$ 中替换得更轻树, 矛盾.
\end{proof}
严格性不能省略. 等权三角形的任意两条边都是最小树, 所以圈上一条并列最重边仍可能属于最小树.
\originalsection{Kruskal 与 Prim}
Kruskal 先把边按权非递减排序, 初始 $A=\varnothing$, 若下一边端点处在当前森林不同分支就加入, 否则跳过. 它不产生圈. 加入时任选其一端所在分支形成割, 所有更轻的跨割边已被扫描; 若未加入, 当时两端已在同分支, 以后也仍在同分支, 不可能跨当前割. 因而当前边为该割的最轻边, 割不穿过 $A$, 安全边定理每步保持可扩为最小树. 最后森林极大, 连通输入使它为生成树, 所以最优. 并查集维护分支, 排序与合并可在 $O(m\log m)$ 时间完成.

Prim 从一点起, 维护已连成的顶点集 $S$ 及其树边 $A$, 每次加入 $S$ 到补集最轻边及其新端点. 割 $\partial S$ 不穿过 $A$, 所以每次安全. 连通性保证到全部顶点前割非空, 最后得到最小生成树. 邻接表与二叉堆实现为 $O((n+m)\log n)$, 矩阵实现为 $O(n^2)$.
\begin{example}
图有边 $ab:1,bc:2,cd:3,ac:4,bd:5,ad:6$. 求最小生成树, 并给出最优性证据.
\end{example}
\begin{solution}
按 Kruskal 顺序加入 $ab,bc,cd$, 已连接四点, 总权6. 非树边 $ac$ 的树路最大权2, $bd$ 的为3, $ad$ 的为3, 都不超过非树边自身权. 这给出最优性证书: 任何其他树的某条非树边加入当前树都不能通过换掉该圈上一边减重. 该证书的充分性将在下段通过固定的基本割逐次交换证明, 避免误用变化中的树路.
\end{solution}
上面的证书充分性应选以原候选树的路径为基准来作交换. 可直接用割证书避免中间树路径变化: 对候选树任一边 $f$, 删它所得基本割上每条非树边 $e$ 的原树路径包含 $f$, 故 $w(f)\le w(e)$. 因而每条树边都是其基本割的一条最轻边. 要把任意其他树改成候选树, 每次添加一个缺失候选边 $f$, 沿其他树中的相应圈删一条跨该基本割且不在候选树中的边 $e$, 权不增加. 有限次变为候选树, 因而原另一树权不小于候选树. 这一论证说明证书如何从局部比较提升到全局最优.
\originalsection{图的交换结构}
\begin{proposition}\label{prop:forestexchange}
若 $A,B$ 都为森林边集且 $|A|<|B|$, 则有 $e\in B\setminus A$ 使 $A\cup\{e\}$ 仍为森林.
\end{proposition}
\begin{proof}
若 $B$ 的每条边都连接 $A$ 的同一分支内部, 对每个 $A$ 分支的顶点集 $V_i$, 森林 $B[V_i]$ 的边数至多 $|V_i|-1$. 求和得到 $|B|\le\sum_i(|V_i|-1)=|A|$, 矛盾. 所以 $B$ 有边跨 $A$ 的两个分支, 加它不成圈.
\end{proof}
这是图拟阵的交换公理. 本册只在图的森林上证明它, 不把一般拟阵优化或多面体理论视为已覆盖. 贪心能正确, 正因为可以把局部选择与其他可行树通过交换联系起来.
\originalsection{随机缩边求最小割}
考虑连通无向多重图, 删除自环, 保留平行边的各自身份. Karger 算法反复在当前全部边中均匀随机选一条, 收缩其端点, 丢弃新自环, 直到只剩两个超级顶点. 剩下平行边形成原图一个割, 算法输出其大小. 收缩可能误删最小割, 因而单次只给概率保证.
\begin{theorem}\label{thm:karger}
在 $n\ge2$ 顶点连通图中, 一次随机缩边保留并输出某个固定最小割的概率至少为 $2/(n(n-1))$.
\end{theorem}
\begin{proof}
固定最小割大小为 $k\ge1$. 条件于先前没有收缩该割的边, 它在当前图仍存在且大小 $k$. 当前图的任何割对应原图的一个割, 所以最小割不小于 $k$, 每个超级顶点度至少 $k$. 若还剩 $r$ 点, 当前边数至少 $rk/2$, 下一步选到固定割边的概率至多 $k/(rk/2)=2/r$. 生存概率至少
\[
\prod_{r=n}^{3}\left(1-\frac2r\right)=\frac2{n(n-1)}.
\]
生存到两点时, 两超级顶点恰为所固定割的两側, 输出该最小割. $n=2$ 时空乘积为1, 与公式一致.
\end{proof}
独立重复 $L$ 次并取最小值, 失败概率至多 $(1-p)^L\le e^{-pL}$, $p=2/(n(n-1))$. 取 $L\ge n(n-1)\log(1/\eta)/2$ 可使失败概率不超过 $\eta$. 每次输出都是真实割, 不会输出低于真正最小值的虚假值. 非连通图的最小割为0, 应先以搜索识别, 不必进行缩边. 平行边不能合并后等概率选择, 否则证明的 $k/m$ 概率已被改变.
\originalsection{习题与解答}
\begin{exercise}\label{ex:mst1}
是否可以用 BFS 树代替最小生成树? 给出一个反例, 即使所有边权非负.
\end{exercise}
\begin{exercise}\label{ex:mst2}
证明任何最小生成树也是瓶颈最小生成树: 它最大边权在全部生成树中最小.
\end{exercise}
习题\ref{ex:mst1}的解答.
\begin{solution}
三角形 $ab:10,ac:10,bc:1$, 以 $a$ 为 BFS 根, 树边为 $ab,ac$, 总权20; 最小树为 $ab,bc$ 或 $ac,bc$, 总权11. BFS 最小化的是根到各点的边数距离, 与生成树总权不同.
\end{solution}
习题\ref{ex:mst2}的解答.
\begin{solution}
取最小树中的最重边 $e$, 删它得到基本割. 若有另一树所有边均严格轻于 $e$, 它必含一条跨割边 $f$, 用 $f$ 换 $e$ 就降低最小树总权, 矛盾. 所以其他树的最大边权至少为 $w(e)$.
\end{solution}

\originalchapter{网络流与割对偶}\label{ch:flow}
\sourceof{18.433 L7--8; 与 Menger、匹配的衔接为编者补充}
\originalsection{容量与守恒}
\begin{definition}
流网络是有限有向多重图, 指定不同顶点源 $s$、汇 $t$, 每条弧 $e$ 有容量 $c(e)\ge0$. 可行流为弧上的实数 $f(e)$, 满足 $0\le f(e)\le c(e)$, 且除 $s,t$ 外每点流出量等于流入量. 流值 $|f|$ 为源的净流出量, 也等于汇的净流入量.
\end{definition}
自环对流值无作用, 可删除. 平行弧与反平行弧各有独立编号. 对 $s\in S,t\notin S$ 的割, 容量为从 $S$ 到补集的弧容量之和, 反向弧不计入容量.
\begin{lemma}[弱对偶]\label{lem:flowweak}
任意可行流与任意 $s$--$t$ 割满足
\[
|f|=\sum_{e:S\to\overline S}f(e)-\sum_{e:\overline S\to S}f(e)\le c(S,\overline S).
\]
\end{lemma}
\begin{proof}
把 $S$ 内各顶点的净流出量求和, 内部弧一进一出抵消, 非源点由守恒贡献0, 所以留下源的流值. 前向流不超过容量, 反向流非负, 得不等式.
\end{proof}
\originalsection{残量网络与增广}
对每条原弧 $e=u\to v$, 建立前向残量弧 $u\to v$ 容量 $c(e)-f(e)$ 和后向残量弧 $v\to u$ 容量 $f(e)$, 只保留正残量弧. 它们保留原边编号与方向身份, 因而反平行原弧不会与另一弧的回退混淆. 前向表示还能增加流量, 后向表示可以撤回已安排流量.

若残量网络有 $s$--$t$ 路, 令瓶颈 $\theta$ 为路上最小残量. 沿前向原弧加 $\theta$, 沿后向减 $\theta$. 简单路不会同时使用同一原弧两种身份. 更新后仍可行, 中间点净变化为0, 源净流出增加 $\theta$. 无后向弧的“只向前加流”容易陷入局部选择, 不能作为一般最大流算法.
\begin{theorem}[最大流最小割]\label{thm:maxflow}
有限网络的最大流值等于最小割容量. 若所有容量为非负整数, 则有整数最大流, 从零流反复沿增广路更新可在有限次后得到它.
\end{theorem}
\begin{proof}
整数容量时, 每次残量和瓶颈都是整数, 一次至少使流值增加1. 流值由源的出弧总容量界住, 所以有限次后无增广路. 令 $S$ 为残量中从源可达点集, 汇不在其中. 每条原前向跨割弧已饱和, 否则还能到达其终点; 每条原反向跨割弧流量为0, 否则其后向残量会从 $S$ 到补集. 所以弱对偶式取等, 流值等于该割容量, 两者都最优.

一般实容量时, 不保证任意选增广路的过程终止. 但可行流构成有限维闭有界集合, 非空, 线性流值在其上达到最大. 最大流若残量仍有路, 正瓶颈会改进流值, 矛盾. 因而它无增广路, 按相同可达集构造割就取等. 这里的存在性仅用有限维紧致性, 不依赖 Menger 或匹配定理.
\end{proof}
\originalsection{Edmonds--Karp 的时间界}
每次在残量网络用 BFS 选择边数最少的增广路, 就得到 Edmonds--Karp 算法. 这不仅选出一条容易找到的路, 还阻止某条弧不断在同一距离层上重复成为瓶颈.
\begin{lemma}\label{lem:residualdistance}
采用最短增广路时, 各顶点从源的残量距离不会减小, 不可达点距离记为无限.
\end{lemma}
\begin{proof}
更新前各旧残量弧 $u\to v$ 满足 $d(v)\le d(u)+1$. 更新后新出现的弧只可能是增广路边的反向. 对增广路上的 $u\to v$, 有 $d(v)=d(u)+1$, 因而新反向弧也满足 $d(u)\le d(v)+1$. 新图中若有比旧距离短的路, 从源沿它逐边应用这些不等式, 该路长仍至少为旧终点距离, 矛盾. 旧时不可达点无法由新反向弧首次到达, 因增广路各点旧时已可达, 所以无限距离也不会下降.
\end{proof}
\begin{theorem}
Edmonds--Karp 至多增广 $O(nm)$ 次, 用邻接表与 BFS 的时间为 $O(nm(n+m))$, 在连通可达网络的通常 $m\ge n-1$ 情形可写为 $O(nm^2)$.
\end{theorem}
\begin{proof}
一次增广至少有一个残量弧饱和, 称为本次临界弧. 若 $u\to v$ 是临界弧, 当时 $d(v)=d(u)+1$. 要使同一身份的残量弧再次出现并再次饱和, 中间必须沿其反向 $v\to u$ 增广以恢复残量. 在那次最短路上, 新距离满足 $d'(u)=d'(v)+1\ge d(v)+1=d(u)+2$, 用了距离不减引理. 所以同一残量弧两次临界之间, 尾点距离至少增加2. 有限可达距离至多 $n-1$, 每种残量身份最多临界 $O(n)$ 次, 总身份数至多 $2m$, 得增广次数. 每次 BFS 扫描 $O(n+m)$ 项, 得时间界. 删除与源汇问题无关的孤立顶点后, 常用界为 $O(nm^2)$.
\end{proof}
证明没有要求容量为整数, 所以算法对精确实数运算也有限终止. 实际计算机处理浮点残量需明确容差; 这里是数学算法时间界, 不把近似浮点比较当作精确证书.
\originalsection{一个流与割的证书}
\begin{example}
网络有弧 $sa:3,sb:2,ab:1,at:2,bt:3$. 求最大流及最小割.
\end{example}
\begin{solution}
先沿 $s,a,t$ 送2, 再沿 $s,b,t$ 送2, 最后沿 $s,a,b,t$ 送1. 得 $f(sa)=3,f(sb)=2,f(ab)=1,f(at)=2,f(bt)=3$, 各中间点守恒, 总流5. 割 $S=\{s\}$ 容量 $3+2=5$, 流值与割容量相等, 所以无需再搜索也能验证最优.
\end{solution}
\begin{center}
\begin{tikzpicture}[x=2.1cm,y=1.15cm]
\node[vertex](s)at(0,0){$s$};\node[vertex](a)at(1,1){$a$};\node[vertex](b)at(1,-1){$b$};\node[vertex](t)at(2,0){$t$};
\draw[arc](s)--node[above left]{$3/3$}(a);\draw[arc](s)--node[below left]{$2/2$}(b);
\draw[arc](a)--node[right]{$1/1$}(b);\draw[arc](a)--node[above right]{$2/2$}(t);\draw[arc](b)--node[below right]{$3/3$}(t);
\end{tikzpicture}
\end{center}
标注为“流量/容量”. 容量上界与守恒是两套约束, 校验证书时二者都必须检查.
\originalsection{匹配、路与对偶的统一}
二部图 $A\sqcup B$ 的最大匹配可化为流: 加源到每个 $A$ 点容量1, 每条二部边作为 $A\to B$ 容量1, 每个 $B$ 点到汇容量1. 整数流使每点最多参与一条正流二部边, 所以流值就是匹配大小; 反过来每个匹配给等值流. 顶点容量拆分则给出第\ref{ch:connectivity}章 Menger 的顶点版本. 同一个守恒与残量机制因此联系了边不相交路、顶点不相交路与匹配.

最大流还可写为线性规划, 最小割为一个特殊的对偶证书. 本册使用残量可达集实际构造这个证书, 不先假设线性规划强对偶. 一般费用流、单纯形、匹配多面体与花算法的复杂度不在本册范围.
\originalsection{习题与解答}
\begin{exercise}\label{ex:flow1}
网络弧为 $sa,sb,ac,ad,bc,ct,dt$, 容量全为1. 若先沿 $s,a,c,t$ 增广, 给出第二次增广路以达到最大流2.
\end{exercise}
\begin{exercise}\label{ex:flow2}
设 $f$ 是一个整数最大流, 说明每个整数 $0\le k\le |f|$ 也能作为某个整数流的值.
\end{exercise}
习题\ref{ex:flow1}的解答.
\begin{solution}
残量路 $s,b,c,a,d,t$ 使用后向弧 $c\to a$, 撤回第一次的 $a\to c$, 改成两条路 $s,b,c,t$ 与 $s,a,d,t$. 流值2, 源总容量2给最优上界. 如果不允许回退, 第一次选择会错误地看似只能送1.
\end{solution}
习题\ref{ex:flow2}的解答.
\begin{solution}
把整数流的圈流先消去, 剩余正流反复沿源汇路抽取一单位, 得 $|f|$ 条单位路流. 取其中任意 $k$ 条相加, 仍满足原容量与守恒, 得整数值 $k$. 这只是可行流的分解, 不要求单位路边不相交, 因为容量可以大于1.
\end{solution}

\Needspace{8\baselineskip}
\originalchapter{矩阵树定理与有向计数}\label{ch:matrixtree}
\sourceof{18.212 L26--32; 行列式展开与最后出边双射补全}
\originalsection{关联矩阵与 Laplace 矩阵}
\begin{definition}
对无向无自环多重图, 临时任意给各边定向. 关联矩阵 $B$ 每列对应一条边 $u\to v$, 在 $u$ 行为1, $v$ 行为 $-1$, 其他为0. 令 $W=\operatorname{diag}(w_e)$, $w_e>0$ 为边权. 加权 Laplace 矩阵 $L=BWB^{\mathsf T}$, 即 $L_{vv}=\sum_{e\ni v}w_e$, $L_{uv}=-\sum_{e:uv}w_e$.
\end{definition}
定向只是计算辅助, 改一列的符号不改变 $BWB^{\mathsf T}$. 对任意向量 $x$, 有
\begin{equation}\label{eq:lapenergy}
x^{\mathsf T}Lx=\sum_{e=uv}w_e(x_u-x_v)^2.
\end{equation}
所以 $L$ 实对称、半正定, 常向量在核中. 核中向量必须沿每条正权边相等, 因而在每个分支上常值. 连通图的核恰为一维, 这既保证电势在加常数后唯一, 也解释为什么树计数要删一行一列.
\originalsection{无向矩阵树定理}
\begin{theorem}[加权矩阵树]\label{thm:matrix-tree}
连通加权图的 $L$ 删去任一顶点 $r$ 对应行列所得矩阵 $L^{(r)}$ 满足
\[
\det L^{(r)}=\sum_{T\text{ 生成树}}\prod_{e\in T}w_e.
\]
无权时各 $w_e=1$, 行列式就是生成树数. 单顶点图的空行列式和空树积均为1.
\end{theorem}
\begin{proof}
删去 $B$ 的 $r$ 行得 $\widetilde B$, 则 $L^{(r)}=\widetilde B W\widetilde B^{\mathsf T}$. Cauchy--Binet 恒等式为: 若 $A$ 为 $k\times m$, $C$ 为 $m\times k$, 则 $\det(AC)=\sum_{|S|=k}\det A_{[:,S]}\det C_{[S,:]}$. 它由行列式的多线性展开得到: 选择重复列索引的项因反对称性相消, 对不同索引的排列求和恰形成两行列式乘积. 用它得到
\[
\det L^{(r)}=\sum_{|S|=n-1}\det(\widetilde B_S)^2\prod_{e\in S}w_e.
\]
若 $S$ 含圈, 圈上各列按方向加减和为0, 所以行列式为0. 若不连通, 取不含 $r$ 的分支, 其顶点行之和在选定列上为零, 同样为0. 若 $S$ 为生成树, 取不为 $r$ 的叶子, 它的行仅有一个 $\pm1$, 沿此行展开, 删除叶子和其边后继续归纳, 得行列式绝对值1. 因而只有生成树项保留.
\end{proof}
\begin{example}
用矩阵树定理求 $K_n$ 和 $K_{a,b}$ 的生成树数.
\end{example}
\begin{solution}
$n\ge2$ 时, $K_n$ 的余子式为 $nI_{n-1}-J_{n-1}$, $J$ 为全1矩阵. 常向量的特征值为1, 与它正交的 $n-2$ 维空间中特征值为 $n$, 得 $n^{n-2}$, 与 Prüfer 计数一致. $n=1$ 时只有空树, 按空行列式约定计数1.

对 $a,b\ge2$, 从大小 $b$ 侧删一点, 余子式为
\[
\begin{pmatrix}bI_a&-J_{a,b-1}\\-J_{b-1,a}&aI_{b-1}\end{pmatrix}.
\]
用第一块消去下左块, 行列式为 $b^a\det(aI_{b-1}-(a/b)J_{b-1})$. 后一矩阵在常向量上的特征值为 $a/b$, 其余 $b-2$ 个特征值为 $a$, 得 $a^{b-1}b^{a-1}$. 若某侧大小1, 图是星, 树数1, 公式也成立.
\end{solution}
\originalsection{有向矩阵树定理}
\begin{definition}
有向加权多重图的出度 Laplace 矩阵满足 $L_{ii}=\sum_{j\ne i}w_{ij}$, $L_{ij}=-w_{ij}$ ($i\ne j$), 其中 $w_{ij}$ 为从 $i$ 到 $j$ 各弧权之和. 指向根 $r$ 的生成树要求每个非根点恰有一条出弧, 根无出弧, 顺着它们最终到达根. 自环不参与树, 因而从矩阵中忽略.
\end{definition}
\begin{theorem}\label{thm:directedmatrix-tree}
删去出度 Laplace 矩阵根行列后的行列式, 等于指向该根的有向生成树权积之和.
\end{theorem}
\begin{proof}
把每个非根行写为出弧贡献之和 $\sum_{e:i\to j}w_e(e_i^{\mathsf T}-e_j^{\mathsf T})$, 删根列时 $e_r$ 视为0. 按行多线性展开, 一项相当于为每个非根点选择一条出弧, 系数为弧权积, 行矩阵由这些差向量组成.

若选择的弧含一个不经过根的有向圈, 圈中行向量之和为0, 行列式为0. 若没有这种圈, 每点沿其唯一出弧必须到达根, 否则有限追踪会陷入圈. 所选弧即指向根的树. 按从叶子向根消元, 或按距根递减排序行列, 该矩阵的行列式为1: 每行在自身列为1, 出弧指向更后位置, 因而为对角全1的三角矩阵. 同时置换行与列不改变行列式符号. 所以恰得到树的权积之和.
\end{proof}
若要计数从根向外的树, 对反图使用 $L^{\mathrm{rev}}_{ii}=\sum_{j\ne i}w_{ji}$, $L^{\mathrm{rev}}_{ij}=-w_{ji}$ ($i\ne j$), 再删根行列. 这同时反转弧和重算对角度数, 不是直接转置原出度矩阵; 单纯转置并不改变余子式行列式. 因此矩阵对角采用入度还是出度、所计树朝向根还是远离根, 必须同时说明, 不能只写“有向版本类似”.
\originalsection{BEST 定理}
\begin{theorem}[BEST]\label{thm:best}
设有限强连通有向多重图各点满足 $d^+(v)=d^-(v)=d(v)\ge1$, 弧有独立标号. 固定从根 $r$ 发出的首弧 $e_0$. 从 $r$ 开始且首弧为 $e_0$ 的 Euler 弧序列数为
\[
t_r\prod_{v\in V}(d(v)-1)!,
\]
其中 $t_r$ 为指向 $r$ 的生成树数. 这里不再按循环旋转或图自同构识别序列.
\end{theorem}
\begin{proof}
从 Euler 序列中, 对每个 $v\ne r$ 取最后一次离开 $v$ 的弧. 这些弧不能成有向圈: 若成圈, 选圈上最后离开时间最晚的点, 它沿所选弧到达下一点后, 下一点仍需离开, 与该点的最后离开时间更早矛盾. 每个非根一点有一个出弧, 无圈保证最终到根, 所以得到指根树.

反向固定这样一棵树. 对每个非根点, 把树弧排在全部出弧次序最后, 其余 $d(v)-1$ 弧任意排列; 在根处把 $e_0$ 排第一, 其余任排. 从根出发, 每次使用当前位置尚未用的第一条出弧. 平衡性保证不能先停在非根点, 因为那里已用入弧会比全部出弧多1, 不可能. 所以先停在根, 所走为闭迹.

若仍有未用弧, 对所有点闭迹已用入出数相等, 未用入出数也相等. 取一个有未用出弧的非根点; 它的树弧排最后, 必还未用. 沿树弧到下一点, 该点有未用入弧, 因而也有未用出弧. 沿树追到根, 推出根有未用出弧, 与停下矛盾. 若原先有未用出弧点就是根, 已直接矛盾. 因而全部弧被使用, 得 Euler 序列. 最后出边与局部次序从序列中唯一恢复, 这是一一对应. 每棵树对应所列阶乘数量, 得公式.
\end{proof}
\begin{example}
两个顶点 $a,b$ 之间各有两条独立标号的反向弧, 从 $a$ 出发固定首弧. 求 Euler 序列数.
\end{example}
\begin{solution}
指向 $a$ 的树需从 $b$ 选一条弧, 有2种, 各点度2, 阶乘积为1, 所以共有2条序列. 固定首弧后, 在 $b$ 首次回 $a$ 的弧有两选择, 其余弧顺序随之确定. 这也核对了根处阶乘因子, 不能误多乘一个 $d(a)$.
\end{solution}
\originalsection{习题与解答}
\begin{exercise}\label{ex:matrix1}
求三角形三条边权为 $a,b,c>0$ 时的生成树权和.
\end{exercise}
\begin{exercise}\label{ex:matrix2}
有向图为三顶点有向圈, 每条弧再复制为两条独立标号平行弧. 求固定根与首弧的 Euler 序列数.
\end{exercise}
习题\ref{ex:matrix1}的解答.
\begin{solution}
删去一个根后的余子式可写为 $\left(\begin{smallmatrix}a+b&-b\\-b&b+c\end{smallmatrix}\right)$, 行列式为 $(a+b)(b+c)-b^2=ab+ac+bc$. 恰好对应三种删去一条边的树权积.
\end{solution}
习题\ref{ex:matrix2}的解答.
\begin{solution}
对两个非根点, 其树弧各有2种选择, 所以 $t_r=4$. 各点出度2, 阶乘积1, 得4. 按局部出弧次序也可验证: 根首弧已固定, 两个非根点各可任选哪条出弧放最后, 共有4种.
\end{solution}

\originalchapter{电网络与随机游走}\label{ch:electrical}
\sourceof{18.212 L30--31; 能量原理与通勤时间恒等式为编者补充}
\originalsection{电势与 Kirchhoff 方程}
本章考虑有限连通无自环无向图, 至少两点. 每边电导 $c_{uv}>0$, 电阻为 $1/c_{uv}$. 平行边可把电导相加, 因而本章用对称权 $c_{uv}=c_{vu}$. 顶点电势为 $\phi_v$, 从 $u$ 到 $v$ 的电流为 $I_{uv}=c_{uv}(\phi_u-\phi_v)$, 反向电流取负. 各点净流出电流组成向量 $b$, 守恒给 $L\phi=b$, 其中 $L$ 为加权 Laplace 矩阵.
\begin{proposition}\label{prop:potential}
若 $\sum_v b_v=0$, 方程 $L\phi=b$ 有解, 解在加常向量的意义下唯一. 固定一点电势为0后唯一.
\end{proposition}
\begin{proof}
第\ref{ch:matrixtree}章已知连通图的 $\ker L$ 为常向量. 实对称矩阵的像为核的正交补, 故和为0的 $b$ 在像中. 两解之差在核中, 即只差常数. 也可直接删根行列求解; 余子式由能量式正定, 因而可逆, 未列的根方程由总和0自动满足.
\end{proof}
若电势在非空边界 $A$ 上固定, 非边界点不注入电流, 则 $\phi_v=\sum_u c_{vu}\phi_u/d_v$, $d_v=\sum_u c_{vu}$. 这是邻点电势的加权平均, 称为调和条件. 固定边界后, 内部主子矩阵正定: 零边界扰动的能量若为0, 沿连通边处处常值并在边界为0, 所以全为0. 因而内部线性系统可逆, 延拓存在. 最大值原理保证解唯一: 两解之差边界为0, 若内部有正最大值, 加权平均迫使所有邻点同样达到它, 沿连通路径传到边界产生矛盾; 负最小值同理.
\originalsection{能量与有效电阻}
\begin{definition}
对不同顶点 $a,b$, 从 $a$ 注入单位电流, 在 $b$ 提取单位电流, 即 $L\phi=\mathbf e_a-\mathbf e_b$. 有效电阻定义为 $R_{ab}=\phi_a-\phi_b$, 与电势加常数无关.
\end{definition}
由能量式, $R_{ab}=\phi^{\mathsf T}L\phi=\sum_{uv}c_{uv}(\phi_u-\phi_v)^2>0$ ($a\ne b$). 严格性因为单位电流要求电势不常值. 边上的电流能量为 $\sum_{uv}I_{uv}^2/c_{uv}$, 与上式相等, 这里每无向边只计一次.
\begin{theorem}[Thomson 原理]\label{thm:thomson}
在所有净流出为 $\mathbf e_a-\mathbf e_b$ 的反对称单位流中, 由电势给出的电流唯一地最小化能量, 最小值为 $R_{ab}$.
\end{theorem}
\begin{proof}
固定每边一个方向, 任意另一单位流写为 $I+J$, 则 $J$ 在每点净流出为0. 能量展开的交叉项为
\[
\sum_{u\to v}\frac{I_{uv}J_{uv}}{c_{uv}}=\sum_{u\to v}(\phi_u-\phi_v)J_{uv}=\sum_u\phi_u\operatorname{div}J(u)=0.
\]
因此 $\mathcal E(I+J)=\mathcal E(I)+\sum J_{uv}^2/c_{uv}\ge\mathcal E(I)$, 等号仅在每边 $J=0$ 时成立.
\end{proof}
对应的 Dirichlet 原理是: 在所有满足 $f(a)=1,f(b)=0$ 的电势函数中, 调和延拓唯一最小化 $\sum c_{uv}(f_u-f_v)^2$, 最小值为有效电导 $1/R_{ab}$. 证明把任意函数写为调和函数加边界为0的扰动, 交叉项为 $h^{\mathsf T}Lf=0$, 因内部 $Lf=0$ 而边界 $h=0$. 电压差为1时的总注入电流 $C$ 满足能量 $C$, 单位电流缩放使 $R=1/C$.
\begin{example}
三角形每边电阻为1. 求两点之间有效电阻, 并用能量检查.
\end{example}
\begin{solution}
对端点电势设为1、0, 第三点由平均关系为 $1/2$. 从电势1点流出的电流为 $1+(1-1/2)=3/2$, 所以有效电导 $3/2$, 电阻 $2/3$. 电势能量为 $1^2+(1/2)^2+(1/2)^2=3/2$, 与总电流一致. 等价地, 直接1欧姆支路与两条串联共2欧姆支路并联, 得 $1/(1+1/2)=2/3$.
\end{solution}
\originalsection{树计数与有效电阻}
记 $\tau=\sum_T\prod_{e\in T}c_e$ 为生成树权和, $F_{ab}$ 为恰有两个分支、$a,b$ 分属不同分支的生成森林权和. 此处权重都是电导, 不是电阻.
\begin{theorem}\label{thm:forestresistance}
对不同顶点 $a,b$, 有 $R_{ab}=F_{ab}/\tau$. 若随机生成树按权积比例选取, 独立标号边 $e=ab$ 被选中的概率为 $c_eR_{ab}$.
\end{theorem}
\begin{proof}
把 $b$ 电势固定为0, 单位电流方程的非根部分为 $L^{(b)}\phi=\mathbf e_a$. 逆矩阵的第 $aa$ 项由余子式给出, 故
\[
R_{ab}=\phi_a=\frac{\det L^{(a,b)}}{\det L^{(b)}}.
\]
分母是矩阵树定理的 $\tau$. 对分子, 删关联矩阵的 $a,b$ 两行并用 Cauchy--Binet, 每项选择 $n-2$ 条边. 含圈时列相关, 行列式为0. 无圈时恰有两个分支; 若某分支不含删去的根, 其顶点行和为0, 仍消失. 只有两根分属两分支的森林保留, 在每棵树中沿非根叶子逐次展开, 行列式绝对值1. 因而分子为 $F_{ab}$. $n=2$ 时分子为空行列式1, 同样正确.

每棵包含指定边 $e$ 的生成树删去 $e$, 唯一得到把其端点分开的两分支生成森林; 反过来给这样的森林加入 $e$ 唯一得到树. 权积乘上 $c_e$, 所以包含 $e$ 的树权和为 $c_eF_{ab}$, 除以总权和得到概率. 平行边以独立编号处理, 双射仍成立.
\end{proof}
该式把确定性的电势差与随机树的边出现概率联系起来. 例如单位三角形每边 $R=2/3$, 每边在三棵等概率生成树中恰出现两次.
\originalsection{随机游走的平稳分布}
\begin{definition}
连通加权图至少两点, 随机游走在 $v$ 处以概率 $P(v,u)=c_{vu}/d_v$ 走到邻点 $u$. 记总加权度 $\mathcal V=\sum_v d_v=2\sum_e c_e$, $\pi(v)=d_v/\mathcal V$.
\end{definition}
\begin{proposition}
$\pi$ 为平稳分布, 且满足细致平衡 $\pi(v)P(v,u)=\pi(u)P(u,v)$.
\end{proposition}
\begin{proof}
两边都等于 $c_{uv}/\mathcal V$. 对 $v$ 求和, 得 $\sum_v\pi(v)P(v,u)=\sum_v c_{vu}/\mathcal V=d_u/\mathcal V=\pi(u)$.
\end{proof}
平稳不等于任意初始分布立即收敛. 二部图上游走每步换侧, 存在周期2, 初始单点分布一般不收敛到 $\pi$. 第\ref{ch:spectral}章采用每步以一半概率原地不动的惰性游走, 才给出统一收敛界.
\originalsection{命中概率与命中时间}
设 $a,b$ 为不同顶点. 从 $v$ 出发先到 $a$ 而后到 $b$ 的概率应准确写为“在第一次到达集合 $\{a,b\}$ 时命中 $a$ 的概率”, 记 $h(v)$. 有 $h(a)=1,h(b)=0$, 内部条件一步得到 $h(v)=\sum_uP(v,u)h(u)$. 因而命中概率恰等于边界电压为1、0的调和电势. 有限连通性保证最终命中: 任一点有长度至多 $n-1$ 的正概率路到边界, 有限点数使在这种时间窗口内命中概率有统一正下界, 尾概率按几何速度下降.

记 $H_{vb}$ 为从 $v$ 首次到 $b$ 的期望步数, $H_{bb}=0$. 同一窗口论证保证期望有限. 一步分解为 $H_{vb}=1+\sum_uP(v,u)H_{ub}$ ($v\ne b$). 记向量 $h^{(b)}_v=H_{vb}$. 乘 $d_v$ 后, $Lh^{(b)}$ 在非目标点为 $d_v$; 因 $L$ 每列和0, 在目标点为 $d_b-\mathcal V$. 所以
\[
Lh^{(b)}=d-\mathcal V\mathbf e_b,
\]
其中 $d$ 为加权度向量.
\begin{theorem}[通勤时间]\label{thm:commute}
$H_{ab}+H_{ba}=\mathcal V R_{ab}$.
\end{theorem}
\begin{proof}
由两个目标对应的方程相减,
$L(h^{(b)}-h^{(a)})=\mathcal V(\mathbf e_a-\mathbf e_b)$.
因此电势向量 $h^{(b)}-h^{(a)}$ 与单位电流电势的 $\mathcal V$ 倍只差常数. 两端相减时常数抵消, 左边为 $(H_{ab}-0)-(0-H_{ba})$, 右边为 $\mathcal V R_{ab}$.
\end{proof}
无权图中 $\mathcal V=2m$. 单位电阻三角形有 $m=3,R=2/3$, 通勤时间为4; 实际每方向期望命中步数为2. 命中时间有方向性, 通勤时间把两方向相加才对应对称电阻.
\originalsection{习题与解答}
\begin{exercise}\label{ex:elec1}
证明增大某些边电导或增加新边不会增大两固定点间有效电阻.
\end{exercise}
\begin{exercise}\label{ex:elec2}
无权路径 $P_n$, $n\ge2$, 的两个端点之间, 求有效电阻与通勤时间.
\end{exercise}
习题\ref{ex:elec1}的解答.
\begin{solution}
使用 Thomson 原理. 旧单位流在新增边取0后仍可行, 对旧边增大电导会使同一流的能量 $I_e^2/c_e$ 不增加. 新网络能量最小值不大于该旧流能量, 因而有效电阻不增加. 无需先求新电势.
\end{solution}
习题\ref{ex:elec2}的解答.
\begin{solution}
每条边是分开两端的桥, 单位流必须在每条边通过1, 能量为 $n-1$, 故电阻 $n-1$. 总度为 $2(n-1)$, 通勤时间为 $2(n-1)^2$. 由左右对称性, 每方向的命中时间均为 $(n-1)^2$.
\end{solution}

\originalchapter{谱方法、扩张与混合}\label{ch:spectral}
\sourceof{18.217 第4章指定谱段落; 正则图限制版与 Cheeger、惰性游走推导为编者补充}
\originalsection{邻接矩阵的谱}
本章邻接谱部分默认顶点数 $n\ge2$. 有限无向简单图的邻接矩阵 $A$ 是实对称矩阵, $A_{uv}=1$ 当且仅当 $uv$ 为边. 由矩阵乘法归纳, $(A^k)_{uv}$ 计数从 $u$ 到 $v$ 的长度 $k$ 行走, 因而 $\operatorname{tr}A^k$ 计数带指定起点的闭行走. 实对称矩阵有一组标准正交特征向量, 全部特征值为实数; 这是本章使用的线性代数前提.

若图为 $d$ 正则, 常向量为特征值 $d$ 的向量. 对任意 $x$, 由 $2|x_ux_v|\le x_u^2+x_v^2$ 得 $|x^{\mathsf T}Ax|\le d\|x\|^2$, 所以全部特征值在 $[-d,d]$ 内. 选常向量为第一特征向量, 其余特征值记 $\lambda_2,\ldots,\lambda_n$, 令 $\lambda_* =\max_{i\ge2}|\lambda_i|$. 图不连通时 $d$ 可能重复, 二部图还可能有 $-d$; 这些都应计入 $\lambda_*$, 不能只看第二大正特征值.
\originalsection{扩张混合引理}
对顶点集 $S,T$, 令 $e(S,T)=\sum_{u\in S,v\in T}A_{uv}$. 这是有序端点计数: 当 $S=T$, 内部每边计两次; 相交集合也依此定义.
\begin{theorem}[扩张混合]\label{thm:mixing}
$n$ 点 $d$ 正则图满足
\[
\left|e(S,T)-\frac{d|S||T|}{n}\right|
\le\lambda_*\sqrt{|S|(1-|S|/n)|T|(1-|T|/n)}.
\]
\end{theorem}
\begin{proof}
把指示向量写为 $\mathbf1_S=(|S|/n)\mathbf1+x$, $x\perp\mathbf1$, 同理 $\mathbf1_T=(|T|/n)\mathbf1+y$. 常向量与正交部分的交叉项为0, 所以差等于 $x^{\mathsf T}Ay$. 在常向量的正交补上, $A$ 的算子范数为 $\lambda_*$, Cauchy--Schwarz 给绝对值不超过 $\lambda_*\|x\|\|y\|$. 直接计算 $\|x\|^2=|S|-|S|^2/n$, 得结论.
\end{proof}
当 $\lambda_*\ll d$, 大集合之间的边数接近按密度均匀分布的预测. 这个结论不是说每一对顶点都接近随机, 而是控制集合总偏差; 小集合仍可能有特殊结构.
\begin{theorem}[Hoffman 独立数界]\label{thm:hoffman}
设 $n\ge2,d>0$, 图为 $d$ 正则, 最小邻接特征值为 $\tau<0$. 则
\[
\alpha(G)\le\frac{-\tau}{d-\tau}n,\qquad \chi(G)\ge1-\frac d\tau.
\]
\end{theorem}
\begin{proof}
由于 $\operatorname{tr}A=0$ 而有正特征值 $d$, 最小值确实为负. 独立集 $S$ 的大小记 $a$, 用上面的分解得 $0=\mathbf1_S^{\mathsf T}A\mathbf1_S=da^2/n+x^{\mathsf T}Ax\ge da^2/n+\tau(a-a^2/n)$. 若 $a>0$, 整理得界; $a=0$ 显然. 每个颜色类独立, 故 $n\le\chi\alpha$, 得染色下界.
\end{proof}
\begin{example}
用最小特征值证明奇圈 $C_5$ 至少需3种颜色.
\end{example}
\begin{solution}
圈的特征向量由五次单位根得到, 特征值为 $2\cos(2\pi j/5)$, $j=0,\ldots,4$. 最小值 $\tau=-(1+\sqrt5)/2$, 所以 $\chi\ge1+4/(1+\sqrt5)=\sqrt5>2$. 染色数为整数, 至少3. 按圈顺次使用颜色 $1,2,1,2,3$ 又给上界3.
\end{solution}
\originalsection{稠密正则图的四步计数}
\begin{proposition}\label{prop:regularquasi}
设 $G_n$ 为 $n$ 点简单正则图, $d_n/n\to p\in(0,1)$. 下列两条件等价: $\lambda_*(G_n)=o(n)$; 带标号四步闭行走数 $\operatorname{tr}A_n^4=p^4n^4+o(n^4)$. 两者均推出对所有 $S,T$ 一致的 $e(S,T)=p|S||T|+o(n^2)$.
\end{proposition}
\begin{proof}
谱展开给 $\operatorname{tr}A^4=d^4+\sum_{i\ge2}\lambda_i^4$, 以及 $\sum_i\lambda_i^2=\operatorname{tr}A^2=nd$. 若 $\lambda_*=o(n)$, 后一求和至多 $\lambda_*^2 nd=o(n^4)$, 且 $d^4=p^4n^4+o(n^4)$. 反之, 从四步计数条件减去 $d^4$, 得非负数和 $\sum_{i\ge2}\lambda_i^4=o(n^4)$, 所以最大项也为 $o(n^4)$, 即 $\lambda_*=o(n)$. 集合偏差由混合引理至多 $\lambda_* n=o(n^2)$, 把 $d/n$ 换为 $p$ 仍产生一致 $o(n^2)$ 误差.
\end{proof}
这里计数允许顶点重复, 等于 $C_4$ 的图同态数; 单射四圈计数仅相差 $O(n^3)$, 但无标号四圈还需除以8. 原课程讨论更一般的不正则拟随机图等价条件; 本命题仅声明正则情形, 不把本证明推广到未控制度序列的图.
\originalsection{归一化 Laplace 矩阵}
现在允许连通无向正电导图, 至少两点, 不含自环. 令 $D=\operatorname{diag}(d_v)$, $L=D-C$, $\mathcal L=D^{-1/2}LD^{-1/2}$. 其特征值按 $0=\mu_1<\mu_2\le\cdots\le\mu_n\le2$ 排列. 下界与核由能量式得到; 上界因为 $\sum_{uv}c_{uv}(f_u-f_v)^2\le2\sum_vd_v f_v^2$. 最小正特征值的变分式为
\[
\mu_2=\min_{f\ne0,\ \sum_vd_vf_v=0}
\frac{\sum_{uv}c_{uv}(f_u-f_v)^2}{\sum_vd_v f_v^2}.
\]
记 $\vol(S)=\sum_{v\in S}d_v$, $c(\partial S)=\sum_{u\in S,v\notin S}c_{uv}$, $\mathcal V=\vol(V)$. 电导扩张常数为
\[
h=\min_{0<\vol(S)\le\mathcal V/2}\frac{c(\partial S)}{\vol(S)}.
\]
\begin{theorem}[Cheeger 不等式]\label{thm:cheeger}
上述连通加权图满足 $\mu_2/2\le h\le\sqrt{2\mu_2}$.
\end{theorem}
\begin{proof}
先取集合 $S$ 的函数 $f=\mathbf1_S-\vol(S)/\mathcal V$, 加权均值为0. 分子为 $c(\partial S)$, 分母为 $\vol(S)(1-\vol(S)/\mathcal V)\ge\vol(S)/2$. 因而 $\mu_2\le2c(\partial S)/\vol(S)$, 对最小集合取界.

反向取 $Lf=\mu_2 Df$ 且加权均值0的特征函数. 正负支撑中至少一侧体积不超过 $\mathcal V/2$, 必有一侧非空; 改符号后令 $g=f_+=\max(f,0)$ 取那侧. 对任意两数, 正部截断满足 $(g_u-g_v)^2\le(f_u-f_v)(g_u-g_v)$. 沿边求和得到
\[
\sum_{uv}c_{uv}(g_u-g_v)^2\le g^{\mathsf T}Lf
=\mu_2\sum_vd_vg_v^2.
\]
每个非空水平集 $S_t=\{v:g_v^2>t\}$ 体积不超过一半. 把每边经过的水平区间长度积分, 有
\[
h\sum_vd_vg_v^2\le\int_0^\infty c(\partial S_t)\,dt
=\sum_{uv}c_{uv}|g_u^2-g_v^2|.
\]
对右边应用 Cauchy--Schwarz, 不超过
\[
\left(\sum_{uv}c_{uv}(g_u-g_v)^2\right)^{1/2}
\left(\sum_{uv}c_{uv}(g_u+g_v)^2\right)^{1/2}
\le\sqrt{2\mu_2}\sum_vd_vg_v^2.
\]
最后一步用了 $(g_u+g_v)^2\le2(g_u^2+g_v^2)$. 非零 $g$ 允许约掉正分母, 得上界.
\end{proof}
证明也给找割的方法: 按特征函数正部数值排序, 检查相邻数值之间的水平集, 其中必有边界比不超过 $\sqrt{2\mu_2}$ 的集合. 无需枚举所有 $2^n$ 个集合.
\originalsection{惰性随机游走的混合}
令 $Q=(I+D^{-1}C)/2$, 每步一半概率停留, 一半按电导走边. 平稳分布仍为 $\pi(v)=d_v/\mathcal V$. 与 $Q$ 相似的实对称矩阵为 $I-\mathcal L/2$, 所以特征值为 $1-\mu_i/2\in[0,1]$; 加入停留消去了负特征值造成的周期振荡.
\begin{theorem}\label{thm:lazymix}
对非负整数 $t$, 从顶点 $s$ 出发, $t$ 步分布 $q_t$ 的全变差距离满足
\[
\frac12\sum_v|q_t(v)-\pi(v)|
\le\frac12\sqrt{\frac1{\pi(s)}-1}\,(1-\mu_2/2)^t.
\]
当 $t=0$ 时, 幂因子按1计.
\end{theorem}
\begin{proof}
分布相对于 $\pi$ 的密度 $a_t(v)=q_t(v)/\pi(v)$ 在加权内积 $\langle f,g\rangle_\pi=\sum_v\pi(v)f(v)g(v)$ 下演化为 $a_{t+1}=Qa_t$, 因细致平衡把转移公式换向. $Q$ 在该内积自伴, 常向量为特征值1. $a_t-1$ 均值0, 范数至多 $(1-\mu_2/2)^t\|a_0-1\|_\pi$, 而初始范数平方为 $1/\pi(s)-1$. Cauchy--Schwarz 给 $\sum_v\pi(v)|a_t(v)-1|\le\|a_t-1\|_\pi$, 即结论.
\end{proof}
扩张、谱隙与混合因此形成同一条链: 小边界产生低能量试探函数, 小谱隙意味着慢衰减; Cheeger 反向保证低能量又能找到小边界. 各公式中的邻接谱与归一化 Laplace 谱须分别辨认.
\originalsection{习题与解答}
\begin{exercise}\label{ex:spec1}
求 $K_n$, $n\ge2$, 的邻接特征值, 用 Hoffman 界恢复独立数与染色数.
\end{exercise}
\begin{exercise}\label{ex:spec2}
在连通、$n\ge2,d>0$ 的无权 $d$ 正则图中, 把 Cheeger 下界写为邻接第二大特征值的形式; 说明为何普通二部游走仍不收敛.
\end{exercise}
习题\ref{ex:spec1}的解答.
\begin{solution}
$A=J-I$, 常向量特征值 $n-1$, 其正交补全为 $-1$. 因而 $\alpha\le n/(n-1+1)=1$, $\chi\ge1+(n-1)=n$, 两者由单顶点集与每点独用颜色取等.
\end{solution}
习题\ref{ex:spec2}的解答.
\begin{solution}
连通正则图有 $\mathcal L=I-A/d$, 故 $\mu_2=1-\lambda_2/d$, 这里 $\lambda_2$ 按代数大小排, 不是绝对值. 所以 $h\ge(d-\lambda_2)/(2d)$. 二部图有邻接特征值 $-d$, 普通转移矩阵有 $-1$, 分布在两侧交替; 即使 $\lambda_2$ 较小也不消除此周期. 惰性矩阵把 $-1$ 变成0.
\end{solution}

\originalchapter{极值图论}\label{ch:extremal}
\sourceof{18.315 L25、L27; 18.217 第2章 Mantel、Turán、二部禁图指定段落}
\originalsection{禁图问题}
给定简单图 $H$, $\operatorname{ex}(n,H)$ 是不含 $H$ 子图的 $n$ 点简单图最大边数. 子图不要求诱导: 多余边可以删除. 极值问题同时需要上界证明与达到或逼近它的构造, 只给一个不等式还未说明答案是否准确.
\begin{theorem}[Mantel]\label{thm:mantel}
$\operatorname{ex}(n,K_3)=\lfloor n^2/4\rfloor$, 平衡完全二部图达到它.
\end{theorem}
\begin{proof}
无三角形图中, 每条边 $uv$ 的邻点集不交, 否则共同邻点构成三角形. 所以 $d(u)+d(v)\le n$. 沿全部边求和, 左边为 $\sum_vd(v)^2$, 右边为 $nm$. Cauchy--Schwarz 给 $\sum_vd(v)^2\ge(2m)^2/n$ ($n>0$), 因而 $m\le n^2/4$; $m=0$ 或 $n=0$ 直接成立. 边数取整数得上界. 大小 $\lfloor n/2\rfloor,\lceil n/2\rceil$ 的完全二部图无三角形, 有乘积 $\lfloor n^2/4\rfloor$ 条边.
\end{proof}
\originalsection{Turán 的精确界}
令 $T_r(n)$ 为把 $n$ 点分成 $r$ 个尽量均衡部分的完全多部图, 允许空部分. 写 $n=qr+s$, $0\le s<r$, 则 $s$ 部大小 $q+1$, 其余大小 $q$, 边数为
\[
t_r(n)=\frac12\bigl(n^2-s(q+1)^2-(r-s)q^2\bigr).
\]
\begin{theorem}[Turán]\label{thm:turan}
对整数 $r\ge1,n\ge0$, 有 $\operatorname{ex}(n,K_{r+1})=t_r(n)$.
\end{theorem}
\begin{proof}
构造 $T_r(n)$ 的团每部至多取一点, 所以不含 $K_{r+1}$, 给下界. 为证明上界, 在禁 $K_{r+1}$ 图中选边数最多者, 再在其中选使孪生类大小平方和最大的图. 两点称孪生, 指它们有相同开邻点集; 不同孪生点必不邻接, 此关系分成若干独立类.

任意非邻接点 $u,v$ 度数相等. 否则删除低度点的全部邻边, 使它成为高度点的孪生, 可增加边数. 操作不产生禁团: 含新点的团换成其孪生点便是原图的同大小团.

若两个不同孪生类 $A,B$ 有一对非邻接点, 两类间所有点均不邻接, 两类公共邻点各自相同, 且两类点度相等. 将 $A$ 全体改成 $B$ 中任一点的孪生, 边数不变, 因 $A$ 独立、原来与 $B$ 无边, 每个被改点仍有原度. 禁团仍不出现, 因新团至多含一个改点, 可换成 $B$ 中点. 原来其他孪生类不会被拆开: 各点对 $A$ 的新邻接关系完全一致. $A,B$ 合成一类, 大小平方和至少增加 $2|A||B|>0$, 与选择矛盾.

所以不同孪生类间全部邻接, 极值图是完全 $k$ 部图. 若 $k>r$, 各选一点产生禁团, 故 $k\le r$. 用空部补成 $r$ 部, 边数为 $(n^2-\sum_i a_i^2)/2$. 若两部大小相差至少2, 从大部移一点到小部, 平方和严格下降, 边数增加. 因而最多边在尽量均衡时取得, 恰为 $t_r(n)$.
\end{proof}
平方和的二级选择使对称化真正结束在极值结构, 避免把“不断复制邻点”当作未经证明会终止的算法. 精确整数式也比直接写 $\lfloor(1-1/r)n^2/2\rfloor$ 更可靠, 后式对一般 $r$ 并非总是准确.
\begin{example}
求10点无 $K_4$ 图的最大边数, 并给达到图.
\end{example}
\begin{solution}
取 $r=3$, $10=3\cdot3+1$, 部分大小4、3、3. 边数为 $4\cdot3+4\cdot3+3\cdot3=33$. 任何无 $K_4$ 图边数至多33, 该完全三部图达到它.
\end{solution}
\originalsection{二部禁图的双计数}
\begin{theorem}[Kővári--Sós--Turán 上界]\label{thm:kst}
设整数 $s,t\ge2$. 若 $n$ 点简单图不含 $K_{s,t}$, 则
\[
m\le\frac12(t-1)^{1/s}n^{2-1/s}+\frac12(s-1)n.
\]
\end{theorem}
\begin{proof}
计数点 $v$ 与其邻点中的 $s$ 元子集 $S$. 同一个 $S$ 的公共邻点至多 $t-1$, 否则得到 $K_{s,t}$; 公共邻点不在 $S$, 因无自环. 所以
\[
\sum_v\binom{d(v)}s\le(t-1)\binom ns\le\frac{t-1}{s!}n^s.
\]
对整数 $d$, $\binom ds\ge(d-s+1)_+^s/s!$, 其中 $x_+=\max(x,0)$. 函数 $(x-s+1)_+^s$ 凸, Jensen 不等式给
\[
\frac n{s!}(2m/n-s+1)_+^s\le\sum_v\binom{d(v)}s.
\]
若平均度至多 $s-1$, 所述界直接成立; 否则取 $s$ 次方根并整理得界. $n=0$ 单独显然.
\end{proof}
取 $s=t=2$, 得无四圈图 $m\le n^{3/2}/2+n/2$. 这个指数与 Turán 的二次边数不同: 禁一个二部图会迫使密度随 $n$ 增大而下降.
\originalsection{无四圈的有限域构造}
取素数 $q$, 在模 $q$ 的域中构造二部图. 一侧为点 $(x,y)\in\mathbb F_q^2$, 另一侧为非竖直线的参数 $(a,b)\in\mathbb F_q^2$, 当 $y=ax+b$ 时邻接. 图有 $2q^2$ 点, 每条线恰过 $q$ 点, 所以有 $q^3$ 边.
\begin{proposition}\label{prop:finitefield}
上述关联图不含四圈, 从而沿 $n=2q^2$ 有 $\operatorname{ex}(n,C_4)\ge n^{3/2}/(2\sqrt2)$.
\end{proposition}
\begin{proof}
四圈意味着两个不同点同时落在两条不同线. 若两点横坐标不同, 两式相减唯一确定斜率 $a=(y_1-y_2)/(x_1-x_2)$, 再唯一确定截距; 域中非零数可除. 若横坐标相同而纵坐标不同, 无同一条非竖直线经过二点. 因而两个不同点至多有一个共同邻线, 排除四圈. 代入顶点数得到边数.
\end{proof}
这是双计数上界的数量级证据, 没有宣称最优常数或每个 $n$ 都直接达到该式. 这里只用素数模运算, 不要求先构造一般素数幂有限域.
\originalsection{习题与解答}
\begin{exercise}\label{ex:ext1}
证明任何 $n$ 点 $m$ 边简单图有一个至少 $m/2$ 条边的二部子图. 能否据此把无三角形等同于二部图?
\end{exercise}
\begin{exercise}\label{ex:ext2}
对 $n\ge1$ 点无 $C_4$ 图直接计数两点的共同邻点, 导出平均度 $\bar d$ 的 $\bar d(\bar d-1)\le n-1$.
\end{exercise}
习题\ref{ex:ext1}的解答.
\begin{solution}
每点独立等概率分到两侧, 每条边跨侧概率 $1/2$, 故跨侧边数期望 $m/2$, 至少一分法不小于期望. 保留跨侧边就是二部子图. 原图本身未必二部: $C_5$ 无三角形但含奇圈, 不能二染色.
\end{solution}
习题\ref{ex:ext2}的解答.
\begin{solution}
任意两点至多一个共同邻点, 所以 $\sum_v\binom{d(v)}2\le\binom n2$. 利用 $\sum_vd(v)^2\ge n\bar d^2$, 左边至少 $n\bar d(\bar d-1)/2$, 得界 ($n>0$). 因而 $m\le n(1+\sqrt{4n-3})/4$, 比直接把一般 KST 式代入略精确.
\end{solution}

\originalchapter{Ramsey、概率方法与偏序}\label{ch:ramsey}
\sourceof{18.315 L3、L22--25、L27 的指定图论主题; 概率证明与偏序算法衔接补全}
\originalsection{Ramsey 数}
\begin{definition}
$R(s,t)$ 为使每张 $n$ 点简单图都有 $s$ 点团或 $t$ 点独立集的最小整数 $n$. 等价地, 给 $K_n$ 每边染红或蓝, 必出现红 $K_s$ 或蓝 $K_t$. 边染色在这里不要求相邻边异色, 与第\ref{ch:edgecolor}章的正常边染色不同.
\end{definition}
\begin{theorem}\label{thm:ramseyrec}
对 $s,t\ge2$, $R(s,t)\le R(s-1,t)+R(s,t-1)$, 因而 $R(s,t)\le\binom{s+t-2}{s-1}$. 边界 $R(1,t)=R(s,1)=1$.
\end{theorem}
\begin{proof}
取 $n=R(s-1,t)+R(s,t-1)$ 并固定点 $v$. 其邻点和非邻点共有 $n-1$ 点, 故邻点至少 $R(s-1,t)$ 或非邻点至少 $R(s,t-1)$, 否则两者和至多 $n-2$. 前一情况在邻点中得到独立 $t$ 集或 $s-1$ 团; 后者加 $v$ 成 $s$ 团. 后一情况得到 $s$ 团或独立 $t-1$ 集; 后者加 $v$ 成独立 $t$ 集. 二项式界由 Pascal 恒等式与边界归纳得到, 同时证明数值有限存在.
\end{proof}
\begin{example}
证明 $R(3,3)=6$.
\end{example}
\begin{solution}
在6点图中, 固定点的5个其他点中至少3个与它邻接, 或至少3个与它不邻接. 前3点若有边, 与固定点成三角形; 无边则为独立3集. 后3点若不全邻接, 一对非邻点与固定点成独立3集; 全邻接则成三角形. 所以6给上界. $C_5$ 及其补图均为五圈, 无三角形, 故 $C_5$ 无3点团或3点独立集, 给下界6.
\end{solution}
\begin{center}
\begin{tikzpicture}[scale=1.1]
\foreach \i in {1,...,5}{\node[vertex](v\i) at({90+72*(\i-1)}:1.25){$\i$};}
\foreach \i/\j in {1/2,2/3,3/4,4/5,5/1}{\draw[edge](v\i)--(v\j);}
\end{tikzpicture}
\end{center}
\originalsection{随机构造 Ramsey 下界}
\begin{theorem}\label{thm:ramseylower}
对整数 $k\ge3$, 有 $R(k,k)>\lfloor2^{k/2}\rfloor$.
\end{theorem}
\begin{proof}
令 $n=\lfloor2^{k/2}\rfloor$, 每条完全图边独立等概率染红或蓝. 一个指定 $k$ 点集为单色的概率是 $2\cdot2^{-\binom k2}$. 单色 $k$ 集总数 $X$ 的期望
\[
\E X=2\binom nk2^{-\binom k2}\le\frac{2^{1+k/2}}{k!}<1.
\]
最后不等式在 $k=3$ 时成立, 此后右边随 $k$ 增大乘 $\sqrt2/(k+1)<1$. 非负整数随机量期望小于1, 必有一次取值为0, 得无单色 $k$ 团的染色. 若 $n<k$, 也直接无 $k$ 集.
\end{proof}
概率证明给存在性, 未必给高效确定构造. 这里先数坏对象而不是计算“所有对象都好”的概率, 因期望可逐项相加, 不要求各坏对象独立.
\originalsection{随机次序的独立集界}
\begin{theorem}[Caro--Wei]\label{thm:carowei}
任意非空简单图满足 $\alpha(G)\ge\sum_v1/(d(v)+1)\ge n/(\bar d+1)$, 其中 $\bar d=2m/n$.
\end{theorem}
\begin{proof}
均匀随机排列顶点, 取在其闭邻点集 $\{v\}\cup N(v)$ 中最早出现的点. 两个相邻点不可能同时入选, 因各自都要求先于对方, 所选集独立. 点 $v$ 入选概率为 $1/(d(v)+1)$, 期望大小为所述求和, 某一次至少达到期望. 对正数应用 Cauchy--Schwarz, $\sum1/(d(v)+1)\ge n^2/\sum(d(v)+1)=n/(\bar d+1)$.
\end{proof}
这一界与最大度贪心的 $n/(\Delta+1)$ 不同, 它能利用度数分布. 对孤立点贡献恰为1, 符合每个最大独立集可以包含它的事实.
\originalsection{大围长与大染色数}
围长为最短圈长度, 森林的围长记为无限. 奇圈阻止二染色, 但去掉所有短圈也不能统一限制染色数.
\begin{theorem}[Erdős 存在性]\label{thm:highgirth}
对任意整数 $g\ge3,r\ge1$, 存在有限简单图围长大于 $g$, 染色数大于 $r$.
\end{theorem}
\begin{proof}
取很大的 $n$, 在 $n$ 点上各边独立以 $p=n^{-1+1/(2g)}$ 出现. 令 $X$ 为长度3到 $g$ 的圈数. 长 $\ell$ 圈候选至多 $n^\ell/(2\ell)$, 所以
\[
\E X\le\sum_{\ell=3}^g\frac{n^\ell p^\ell}{2\ell}=O_g(n^{1/2}).
\]
Markov 不等式给 $\Prob(X\ge n/4)=O_g(n^{-1/2})\to0$.

令 $t=\lceil n/(2r)\rceil$, 任一指定 $t$ 集独立的概率为 $(1-p)^{\binom t2}\le\exp(-p\binom t2)$. 存在独立 $t$ 集的概率由联合界至多 $2^n\exp(-p\binom t2)$, 趋于0, 因指数中负项数量级为 $n^{1+1/(2g)}$, 超过正项 $n\log2$. 所以足够大 $n$ 时可同时有 $X<n/4$ 且 $\alpha<t$.

在这张图中, 对每个短圈挑一个顶点删除, 共删少于 $n/4$ 点. 新图不出现新圈, 原短圈均被打破, 围长大于 $g$, 保留 $N>3n/4$ 点且独立数至多 $t-1<n/(2r)$. 任何染色的每类都是独立集, 故 $\chi\ge N/\alpha>3r/2>r$.
\end{proof}
删去坏对象的少量支撑而保留全局性质, 称为修改法. 这里删点保持独立数上界; 若只删边, 可能增加独立集, 不能直接沿用结论.
\originalsection{局部引理与超图二染色}
超图的每条边可包含多个顶点. 正常顶点二染色要求每条超边都有两种颜色, 并不要求超边内部两两异色.
\begin{theorem}[对称局部引理]\label{thm:lll}
设坏事件 $A_1,\ldots,A_N$ 各概率至多 $p$. 有依赖图最大度 $D$, 使每个事件与所有非邻事件共同生成的事件族独立. 若某个 $x\in(0,1)$ 满足 $p\le x(1-x)^D$, 则 $\Prob(\bigcap_i\overline A_i)>0$. 特别地, $D\ge1$ 且 $ep(D+1)\le1$ 时成立.
\end{theorem}
\begin{proof}
归纳证明对不含 $i$ 的任意指标集 $S$, 条件概率 $\Prob(A_i\mid\bigcap_{j\in S}\overline A_j)\le x$, 且所用条件事件有正概率. $S$ 为空时由 $p\le x$ 成立. 假设较小集合结论成立, 先按条件概率连乘, 避开 $S$ 的概率至少 $(1-x)^{|S|}>0$.

把 $S$ 分为 $i$ 的邻事件指标 $R$ 与非邻指标 $Q$. 在已避开 $Q$ 的条件下, 分子 $\Prob(A_i\cap\bigcap_{j\in R}\overline A_j\mid\bigcap_{j\in Q}\overline A_j)$ 至多 $\Prob(A_i)=p$, 用了联合独立性. 分母为避开 $R$ 的条件概率; 依次加入 $R$ 指标, 每一步所条件的指标数小于 $|S|$, 由归纳至少 $1-x$. 所以所求概率至多 $p/(1-x)^{|R|}\le p/(1-x)^D\le x$. 全部事件按次序避开, 总概率至少 $(1-x)^N>0$.

最后取 $x=1/(D+1)$. 因 $(D/(D+1))^D\ge e^{-1}$, 有 $x(1-x)^D\ge1/(e(D+1))$. $D=0$ 时事件联合独立, $p<1$ 即保证全部避开概率正.
\end{proof}
\begin{corollary}\label{cor:hypergraph}
有限 $k$ 均匀超图中, 若每条超边至多与 $D$ 条其他超边相交, 且 $e(D+1)\le2^{k-1}$, 则存在正常顶点二染色.
\end{corollary}
\begin{proof}
每点独立等概率染两色, 每条边单色概率 $2^{1-k}$. 不相交超边的坏事件依赖于互不相交的独立顶点变量, 满足所需联合独立性. 应用局部引理, 可避开全部单色超边. 若 $D=0$, 用独立情形即可.
\end{proof}
只说每两事件独立不够; 局部引理需要事件与整组非邻事件的联合独立. 模型中底层独立变量及其相交关系明确给出了这个条件.
\originalsection{偏序的链与反链}
有限偏序集的链是两两可比较的子集, 反链是两两不可比较的子集. 最长反链大小记 $a$, 最长链大小记 $h$.
\begin{theorem}[Dilworth]\label{thm:dilworth}
有限偏序集可分成 $a$ 条链, 且更少的链不能覆盖全部元素.
\end{theorem}
\begin{proof}
为每个元素造左右两个副本 $x_L,x_R$, 当 $x<y$ 时加边 $x_Ly_R$. 最大匹配大小记 $\nu$. 匹配边在原元素上构成每点入出度至多1的有向链, 不能成圈, 因偏序严格增加. 共 $n-\nu$ 条链, 覆盖全部点. 反过来任一 $k$ 链划分, 给每条链的相邻元素加匹配边, 共 $n-k$ 条, 所以最少链数为 $n-\nu$.

由 König 定理, 该二部图有大小 $\nu$ 的顶点覆盖 $C$. 取左右副本都不在 $C$ 的原元素组成 $U$. 被 $C$ 排除的原元素至多 $\nu$ 个, 故 $|U|\ge n-\nu$. $U$ 内若有 $x<y$, 边 $x_Ly_R$ 未被覆盖, 矛盾, 所以 $U$ 是反链. 任一反链与一条链至多交一点, 因而 $a\le n-\nu\le|U|\le a$, 全部取等.
\end{proof}
\begin{theorem}[Mirsky]\label{thm:mirsky}
有限偏序集可分成 $h$ 个反链, 且更少的反链不能覆盖全部元素.
\end{theorem}
\begin{proof}
给 $x$ 赋以 $x$ 为最大元素的最长链长度 $\ell(x)$. 若 $x<y$, 可在该链后添 $y$, 所以 $\ell(y)>\ell(x)$. 同一数值层无可比较点, 是反链; 数值在1到 $h$ 间, 得至多 $h$ 层. 最长链与每个反链至多交一点, 至少需 $h$ 层. 空集有 $h=a=0$, 结论按空划分成立.
\end{proof}
\begin{corollary}[Erdős--Szekeres]\label{cor:es}
任意 $rs+1$ 个互不相同实数的序列, 有长 $r+1$ 的递增子序列或长 $s+1$ 的递减子序列, $r,s\ge1$.
\end{corollary}
\begin{proof}
给第 $i$ 项记两个数: 以它结尾的最长递增长度 $p_i$, 最长递减长度 $q_i$. 若两种目标均没有, 每对 $(p_i,q_i)$ 在 $\{1,\ldots,r\}\times\{1,\ldots,s\}$ 内. 对 $i<j$, 若第 $i$ 项较小则 $p_j\ge p_i+1$, 较大则 $q_j\ge q_i+1$. 所以数对全部不同, 但有 $rs+1$ 对而仅 $rs$ 个位置, 矛盾.
\end{proof}
\begin{theorem}[Sperner]\label{thm:sperner}
集合 $\{1,\ldots,n\}$ 的子集组成包含关系偏序. 任一反链 $\mathcal A$ 满足
\[
\sum_{S\in\mathcal A}\frac1{\binom n{|S|}}\le1,\qquad
|\mathcal A|\le\binom n{\lfloor n/2\rfloor}.
\]
全部中间大小的子集达到第二界.
\end{theorem}
\begin{proof}
每个顶点排列给一条从空集到全集的链, 依次加入排列元素. 反链与它至多交一次. 对固定 $k$ 元集 $S$, 恰有 $k!(n-k)!$ 个排列在该链上经过 $S$. 对排列与所遇反链元素双计数, 得 $\sum_{S\in\mathcal A}|S|!(n-|S|)!\le n!$, 即第一式. 相邻二项式系数比为 $(n-k)/(k+1)$, 所以最大值在中间. 每个求和项至少为该最大系数的倒数, 得第二式. 同大小不同子集互不包含, 中间层确为反链. $n=0$ 时唯一空集也符合.
\end{proof}
\originalsection{连通性保证 Hamilton 圈}
\begin{theorem}[Chvátal--Erdős]\label{thm:chvatalerdos}
至少3点的简单图若顶点连通度 $\kappa(G)\ge\alpha(G)$, 则有 Hamilton 圈.
\end{theorem}
\begin{proof}
若 $\alpha=1$, 图完全, 结论显然. 否则 $k=\kappa\ge\alpha\ge2$. 最小度至少 $k$, 图中有长至少 $k+1$ 的圈: 取最长路, 末点所有邻点都在路中, 至少 $k$ 个邻点保证它与最早的邻点围成至少 $k+1$ 的圈. 取最长圈 $C$, 若未覆盖全部点, 取圈外点 $v$.

存在从 $v$ 到 $C$ 的 $k$ 条内部不相交路, 各路在不同的 $x_1,\ldots,x_k$ 首次进入 $C$. 这一扇形形式由 Menger 拆点法得出: 沿用第\ref{ch:connectivity}章的拆点网络, 源点 $v$ 不拆, 圈上点也容量1, 圈外非源点容量1, 圈上点接新汇, 原边对应弧与接汇弧容量均取 $n+1$; 若少于 $k$ 条路, 相应不到 $k$ 个点的割会把 $v$ 与剩余非空 $C$ 分开, 违反 $k$ 连通. 路可在首次入圈处截断.

沿同一方向记 $x_i^+$ 为圈上后继. $v$ 不邻接任何 $x_i^+$, 否则用扇路从 $x_i$ 经 $v$ 到 $x_i^+$ 替换圈的一条边, 得更长圈. 不同后继 $x_i^+,x_j^+$ 也不邻接: 用两条扇路连 $x_i,x_j$, 加该后继边, 再沿原圈两段反向和正向连接, 得覆盖全部原圈点及圈外点的更长圈. 具体是从 $x_i$ 经圈外路到 $x_j$, 沿圈反向到 $x_i^+$, 跨到 $x_j^+$, 再沿圈正向回 $x_i$. 各段内部不交. 因而 $\{v,x_1^+,\ldots,x_k^+\}$ 为大小 $k+1$ 独立集, 与 $\alpha\le k$ 矛盾. 最长圈必须覆盖全部点.
\end{proof}
这与 Dirac 度数条件互补: 这里比较全局连通度与独立数, 并用内部不相交路把局部插点失败转化成过大的独立集.
\originalsection{递增边标号的长迹}
\begin{proposition}[Graham--Kleitman 的平均度界]\label{prop:increasingtrail}
简单图 $n\ge1$ 的每条边赋不同实数标号, 则有一条沿边标号严格递增的迹, 长度至少 $2m/n$. 迹可重复顶点, 不能重复边; 完全图因而有长至少 $n-1$ 的递增迹.
\end{proposition}
\begin{proof}
按标号递增处理边, 初始每点记录以该点结尾的最长递增迹长度0. 处理 $uv$ 前记录为 $a,b$, 同时更新为 $a'=\max(a,b+1)$, $b'=\max(b,a+1)$. 用当前新边延长另一端已有迹, 合法, 因旧迹边标号都更小; 两次更新必须使用旧值. 反过来以当前边结束的递增迹之前只能使用旧边, 所以记录仍是确切最长长度.

每次两端长度总增至少2: $a=b$ 时各增1; 若 $a\ge b+1$, $a$ 不变而 $b$ 升到 $a+1$, 增量至少2; 另一情况对称. 最后所有记录和至少 $2m$, 故某个记录至少平均数 $2m/n$, 对应一条所求迹. 这并未证明同样长的递增简单路, 重复顶点在此允许.
\end{proof}
\originalsection{习题与解答}
\begin{exercise}\label{ex:prob1}
对5点路径的自然偏序 $1<2<3<4<5$, 和仅有关系 $1<3,2<3,2<4$ 的4点偏序, 分别求最少链划分数及反链划分数.
\end{exercise}
\begin{exercise}\label{ex:prob2}
若 $k$ 均匀超图共有 $M$ 条边, 证明 $M<2^{k-1}$ 足以二染色. 比较这一条件与局部引理的用途.
\end{exercise}
习题\ref{ex:prob1}的解答.
\begin{solution}
全序为一条链, 最长反链1, 链划分数1; 最长链5, 反链划分数5. 第二个偏序分成链 $\{1,3\},\{2,4\}$, 最少链数2, 因 $\{1,2\}$ 为反链. 最长链2, 反链层可取 $\{1,2\},\{3,4\}$, 最少反链划分数2. 这里列出的关系已经传递封闭, 没有隐含更长链.
\end{solution}
习题\ref{ex:prob2}的解答.
\begin{solution}
随机二染色下单色边数期望为 $M2^{1-k}<1$, 故有0条坏边的染色. 条件控制总边数; 局部引理控制每条边相交的其他边数, 即使 $M$ 很大也可能适用. 两个方法分别利用坏对象少和依赖稀疏, 不能把局部度 $D$ 错换为无条件的全图顶点度.
\end{solution}

\originalchapter{来源对应与阅读边界}\label{ch:sources}
\originalsection{主源的选择}
比较以官方公开材料为依据. 6.042J Spring 2015 有完整920页教材, 但逻辑、一般计数与概率占大量篇幅, 因此只补图论基础. 18.314 Fall 2014 有图论教学进度和少数矩阵树讲义, 主阅读依赖商业教材, 不适合作为完整公开主源. 18.315 Spring 2005 的讲义目录提供染色、多项式、Hamilton、Menger、Hall、偏序与极值的专门主线, 因此采用其19份图论相关记录. 18.217 Fall 2019 偏研究级极值与加法组合, 仅用指定极值、谱段落. 更新的 18.225 Fall 2023 作者书稿344页, 独立版权与图片许可说明存在, 此次只保留官方链接及核查元数据.

18.315 官方课程标签为 Spring 2005, 部分手写页却标9--11月2005的课堂日期. 本稿保留官方标签, 并记录原件日期差异, 不自行改成另一学期. PDF 创建、修改时间是数字文件加工时间, 不能当作授课年份. 18.315 记录者 Amanda Redlich、18.212 记录者 Andrew Lin、18.217 班级学生的归属均按官方目录, 不把 PDF 制作者字段误当作者.
\originalsection{18.315 逐文件对应}
下表页数均为 PDF 物理页数. 采用指纳入本书主线; 补证表示源只给提纲或改用自足证明. 归档原件并不表示其中每句话都已翻译.
\begin{longtable}{p{13mm}p{12mm}p{130mm}}
\toprule
文件 & 页数 & 采用范围与边界\\\midrule\endhead
L1 & 2 & Ramsey 的有限性与递推; 多面体刚性讨论不收入.\\
L3 & 2 & 超图二染色的概率思路; 加法数论和凸点集版本不收入.\\
L4 & 2 & 染色重配置与随机二染色; 矩形三染色高度函数不收入.\\
L5 & 2 & 贪心、Brooks; 完整证明补齐; 特定相交超图界不收入.\\
L6 & 2 & 五色、Vizing; 两者补成完整证明.\\
L7 & 2 & 二部图边染色、Heawood 上界; 不宣称证明所有曲面精确染色数.\\
L8 & 1 & 重配置及边数、独立数的染色界; Glauber 速度与二次直径界不收入.\\
L9 & 1 & 染色多项式、包含排除、断圈定理.\\
L10 & 2 & 无圈定向、Stanley 取值、Tutte 两种定义.\\
L11 & 2 & Tutte 特殊值和活动展开; 最小、最大边约定明确区分.\\
L12 & 2 & 圈的多项式与生成树计数联系; Gessel 完全图活动公式不收入.\\
L19 & 1 & Hamilton 必要条件; 特定长路径平面图界不收入.\\
L20 & 2 & Grinberg、Dirac; 顶点传递图猜想不作已证结果.\\
L21 & 2 & 奇度图 Hamilton 奇偶论证、对称群相邻交换构造; 四连通平面图定理仅提及.\\
L22 & 2 & 路分离思想与 Chvátal--Erdős; 一般群的三个对合构造不收入.\\
L23 & 2 & Menger 核心及流的衔接; Gallai--Milgram 路径覆盖不收入.\\
L24 & 2 & Hall、Dilworth、Mirsky、单调子序列; 证明依赖重新排列.\\
L25 & 2 & Sperner、Mantel、递增迹平均度界.\\
L27 & 1 & Turán 精确界与随机次序独立数界.\\\bottomrule
\end{longtable}
L2、L13--18、L26、L28--39不是本次原件冻结选择; 它们含加法组合、纽结、集合系统、排列模式、分拆、铺砌等其他主题. 未选原件的官方链接仍可从课程目录查到, 不将未读资料冒称纳入正文.
\originalsection{补充资料的逐文件对应}
\begin{longtable}{p{33mm}p{12mm}p{110mm}}
\toprule
课程及文件 & 页数 & 正文采用范围\\\midrule\endhead
18.433 L1--3 & 9 & 匹配、增广路; 一般匹配的花算法不收入, Tutte 结构证明另行补齐.\\
18.433 L7--8 & 7 & 流、割、增广与距离不变量; 算法条件及多重弧身份补全.\\
18.433 L9 & 3 & 最小割及随机收缩; 用概率生存界给完整推导.\\
18.212 L22 & 5 & 仅末两页的 Cayley 树计数引入; 前面的分拆恒等式不收入.\\
18.212 L23 & 6 & Cayley、Prüfer 与度数出现次数; 替代多项式证明不逐项翻译.\\
18.212 L26 & 5 & Laplace 矩阵、矩阵树及经典例子.\\
18.212 L27 & 5 & 矩阵树的相关讨论; 不展开全部行列式计算变体.\\
18.212 L28 & 5 & 关联矩阵、Cauchy--Binet、树子式; 权重补充.\\
18.212 L29 & 5 & 有向矩阵树的命题; 用行多线性给完整替代证明.\\
18.212 L30 & 4 & 有向树权与电网络引入; 另一抵消证明不重复.\\
18.212 L31 & 5 & Kirchhoff、电阻与二根森林比值、命中概率.\\
18.212 L32 & 4 & 有向 Euler 与 BEST; 首弧和计数等价约定补齐.\\
18.217 第2章 & 26 & PDF第1--4页 Mantel、Turán; 第6--8页 KST; 第12页起的代数构造思想仅采用无四圈关联图. 其余研究级禁图内容不收入.\\
18.217 第4章 & 19 & PDF第1--6页的稠密谱思路, 第6--8页的扩张混合. 拟随机条件仅证明正则限制版; Cayley 拟随机、Alon--Boppana、Ramanujan、稀疏正则性不收入.\\
6.042 开放教材 & 920 & 第9章有向图与调度, 第11章图、染色、连通、树, 第12章初等平面图. 稳定婚姻及非图论章节不收入. 原件完整留存以保持许可和出处.\\\bottomrule
\end{longtable}
6.042 主要位置对应: 印刷页323--344为 PDF第332--353页的有向图指定理论段落; 印刷页393--433为 PDF第402--442页的简单图主干; 印刷页473--492为 PDF第482--501页的平面图理论. 原件题目分散在章节后部, 本书练习均按本稿重新表述, 未整批搬运其题集.
\originalsection{编者补充与复核记录}
Havel--Hakimi 的换边、完整 Tutte 存在性证明、Brooks 的结构引理、Vizing 扇形、最大流与 Menger 的非循环衔接, 以及搜索、最短路、最小生成树的算法不变量, 是使主线自足的重构或补证. 电网络的能量与通勤时间、Hoffman、Cheeger、惰性游走、稀疏随机构造及局部引理也明确属于补充. 经典数学定理的名称表示归属, 不表示照搬外部商业教材的文字或图形.

原件共34份、1062页, 按 SHA-256 无完全相同文件; 数学内容重复在正文合并, 不用字节哈希冒充语义去重. 来源清单逐项给原件链接、作者年份、页数、文件哈希、提取方式与许可判断. 扫描手写稿的 OCR 仅作检索, 读数学时使用原图. 实际版权提示与 OCR 将数学符号误识为版权符号的情形分别记录. 作者书稿和外部商业教材只留链接, 不收其受限正文、图片或题面.

独立审读记录、逐页视觉记录、原生 XeLaTeX 编译、目录跳转核验、源码 ZIP 干净重编和远端字节校验分别保存. 自动排版检查不能证明数学正确, 小规模计算不能代替一般证明, 独立代理审读也不冒称人类专家审定. 发布所绑定的是最终稿实际哈希.

\begin{thebibliography}{9}
\bibitem{pak} Igor Pak (授课), Amanda Redlich (记录). \textit{18.315 Combinatorial Theory, Spring 2005}. MIT OCW. 图论主干19份课堂记录; \href{https://ocw.mit.edu/courses/18-315-combinatorial-theory-introduction-to-graph-theory-extremal-and-enumerative-combinatorics-spring-2005/pages/lecture-notes/}{官方目录}.
\bibitem{mcs} Eric Lehman, F. Thomson Leighton, Albert R. Meyer. \textit{Mathematics for Computer Science}. 修订版2015-05-18, 6.042J Spring 2015, MIT OCW. 指定使用第9、11、12章相关部分; \href{https://ocw.mit.edu/courses/6-042j-mathematics-for-computer-science-spring-2015/resources/mit6_042js15_textbook/}{原件}.
\bibitem{vempala} Santosh Vempala. \textit{18.433 Combinatorial Optimization, Fall 2003}. MIT OCW. L1--3、L7--8、L9; \href{https://ocw.mit.edu/courses/18-433-combinatorial-optimization-fall-2003/pages/lecture-notes/}{官方目录}.
\bibitem{postnikov} Alexander Postnikov (授课), Andrew Lin (记录). \textit{18.212 Algebraic Combinatorics, Spring 2019}. MIT OCW. L22--23、L26--32指定图论内容; \href{https://ocw.mit.edu/courses/18-212-algebraic-combinatorics-spring-2019/pages/lecture-notes/}{官方目录}.
\bibitem{zhao} Yufei Zhao (授课及编辑), MIT 18.217 Fall 2019 班级学生 (记录). \textit{Graph Theory and Additive Combinatorics}. MIT OCW. 选第2章极值基础及第4章谱工具; \href{https://ocw.mit.edu/courses/18-217-graph-theory-and-additive-combinatorics-fall-2019/pages/lecture-notes/}{官方目录}.
\bibitem{terms} MIT OpenCourseWare. \href{https://ocw.mit.edu/pages/privacy-and-terms-of-use/}{使用条款}. 原件单独许可优先; 改编许可 \href{https://creativecommons.org/licenses/by-nc-sa/4.0/}{CC BY-NC-SA 4.0}.
\end{thebibliography}
\end{document}
